% p2_results.tex
%
% Self-contained write-up of the characteristic-2 results of the research
% session on the Krajicek pseudo-solution pipeline (arXiv:2609.35927) for
% -PHP_n.
%
% Attribution: the pipeline, the solution condition, and the chi-task are due
% to J. Krajicek (arXiv:2609.35927).  All other theorems, propositions, and
% lemmas below are original to this research session (2026-10-03) and are not
% peer-reviewed.  This is a draft for the corpus owner; it has not been
% submitted anywhere.
%
% Revised 2026-10-03 to the corpus's post-correction state: retracts the
% two-phase exact law (simulator artifact) and the all-degrees non-adaptive
% cap; adds the exact certification floor (Theorem F), the rc parity tree,
% and the budget reconciliation; measured points are split into
% hybrid-artifact and literal rows.
%
% Revised 2026-10-03 (final pass) to the corpus's post-kernel state: integrates
% the kernel-structure findings (one design space, V(n,d)^rho = V(2d,d); the
% channel law restated with parity-locked full rows; the i.i.d. law resting on
% the zeroing completion after the disjoint-support route fails at d >= 2),
% the adjacency certificates (Lemma C) with the confirmation-chain
% measurements, the Theorem-F-under-parity-locking remark, and the
% theorem-faithfulness table; supersedes the design-space readings discussion
% and the canonical reversed-certificate scenario.
%
% Compile: tectonic p2_results.tex  (or pdflatex, twice, for cross-references).

\documentclass[11pt]{article}

\usepackage{amsmath}
\usepackage{amsthm}
\usepackage{amssymb}
\usepackage[margin=1.1in]{geometry}

\newtheorem{theorem}{Theorem}[section]
\newtheorem{proposition}[theorem]{Proposition}
\newtheorem{lemma}[theorem]{Lemma}
\newtheorem{corollary}[theorem]{Corollary}

\theoremstyle{definition}
\newtheorem{definition}[theorem]{Definition}

\theoremstyle{remark}
\newtheorem{remark}[theorem]{Remark}
\newtheorem{openproblem}{Open Problem}

\newcommand{\F}{\mathbb{F}}
\newcommand{\PHPn}{\neg\mathrm{PHP}_n}
\newcommand{\Des}{\mathrm{Des}}
\newcommand{\Sp}{\mathrm{Span}}
\newcommand{\Prob}{\mathbb{P}}
\newcommand{\AC}{\mathrm{AC}}
\newcommand{\lab}{\mathrm{lab}}

\title{Adaptive decision-tree caps and the exact certification floor
for the pseudo-solution pipeline of $\PHPn$ at characteristic $2$:\\
the answer channel, depletion bounds, the AND-query no-lift, the
certification floor, the adjacency certificates, and the corrected
two-phase tree}

\author{Research session of October 3, 2026\thanks{
Draft prepared for the corpus owner; not submitted and not peer-reviewed.
The pipeline $\Omega(n,d)$, the solution condition
(Definition~\ref{def:solution}), and the chi-task
(Theorem~\ref{thm:chi}) are due to J.~Kraj\'i\v{c}ek \cite{Kra26}, building
on \cite{BIKPPS97}; they are quoted here from the printed definitions of
\cite{Kra26}.  All remaining results
(Theorems~\ref{thm:channel}, \ref{thm:B}, \ref{thm:and},
\ref{thm:twophase}, \ref{thm:floor}, Propositions~\ref{prop:fixed},
\ref{prop:nonadaptsingle}, \ref{prop:nonadaptall}, \ref{prop:rigidity},
\ref{prop:rcparity}, and
Lemmas~\ref{lem:degonespan}, \ref{lem:completion}, \ref{lem:factor},
\ref{lem:depletion}, \ref{lem:linear}, \ref{lem:certificates},
\ref{lem:rescue}, \ref{lem:classb}, \ref{lem:adjacency}) are original
to this research session.  This revision retracts the earlier exact
two-phase law and
its optimality conjecture (a simulator artifact;
Remark~\ref{rem:artifact}) and the all-degrees non-adaptive cap
(Remark~\ref{rem:retraction}); the final revision integrates the
kernel-level verification (one design space, parity-locked full rows),
the adjacency certificates (Lemma~\ref{lem:adjacency}) and the
confirmation-chain measurements, the Theorem-\ref{thm:floor}
parity-locking remark, and the theorem-faithfulness table
(Table~\ref{tab:faithfulness}).}}

\date{October 3, 2026}

\begin{document}
\maketitle

\begin{abstract}
We study the pseudo-solution pipeline $\Omega(n,d)$ of Kraj\'i\v{c}ek
\cite{Kra26} for the negated pigeonhole principle over $\F_2$, the
reduction target of the $\AC[2]$-Frege lower-bound problem.  We prove an
exact answer-channel law: killed pairs answer deterministically, free
pairs answer fair bits marginally and jointly on every row-incomplete
query set, a fully queried free row is \emph{parity-locked} (answer XOR
$1$; only $2^{2d-1}$ of the $2^{2d}$ patterns occur), and the row sums
answer the determined $0$ on every pigeon.  A kernel-level verification
proves $V(n,d)^{\rho}=V(2d,d)$ exactly: one design space, the
outer/canonical ambiguity dissolves, the printed vanishing condition
itself forces the row parities, and the i.i.d.\ law rests on the
zeroing completion (the earlier disjoint-support route fails at
$d\ge2$).  On this law we prove a hypergeometric depletion bound via
negative association; exact success caps for fixed-label trees, for
non-adaptive variable- and monomial-class trees, and for all adaptive
single-variable trees within their budget; and a no-lift theorem: the
per-hit freeness posterior of a degree-$2$ AND query equals, by an
algebraic identity, the single-variable Bayes posterior.  In the coin
model, our main theorem settles the unbounded-budget certification
problem at degree $1$: every adaptive degree-$1$ tree errs with
probability at least
$\mathrm{err}^*=(1-2d/n)\,(2d+1)(2d)!\,/\,2^{(2d+1)2d}
=2^{-\Theta(d^2)}$, attained by a non-adaptive full-scan counting
tree, through sound line-count certificates and a near-permutation
obstruction; under the pipeline's parity-locking the order survives
and the exact constant is open.  An earlier draft asserted an exact
law $1-2^{-(2d+1)}$ for the two-phase certification tree; we retract
it as a simulator artifact (the simulator certified freeness on a
row-XOR condition that no degree-$1$ query implements), and against
the literal pipeline the tree is dead outright.  The per-hit caps are
sharp but not per-tree: two adjacent $1$s certify a free pair
\emph{with certainty} (a matched pair forces its row-neighbors to
$0$), and a scan-then-confirm chain, sound on the literal pipeline,
measures success $0.9317$ at budget $1000$ at both $(32,2)$ and
$(64,4)$ ($0.9417$ on the exact pipeline at $(32,2)$) with posterior
$1.0000$ on every certified output, the certificate rate following
$1-e^{-B/400}$; the scaling of that constant in $(n,d)$ is the
sharpest open problem.  We tabulate which results are kernel-faithful
(the caps, the posteriors, the completion, both surviving certificate
mechanisms) and which are coin-model-only (the retracted two-phase
law, the floor's constant, the AND identity's pipeline transfer).
Every attacker remains over the printed budget $e'=d\log k$: the open
problem is the \emph{budgeted} certification floor over the
certificate trichotomy --- parity (dead), counting (inert within
budget), adjacency (alive) --- with the chi-hypothesis open in exactly
its printed budgeted form, together with the odd-characteristic
analogue.
\end{abstract}

\section{Introduction}\label{sec:intro}

\subsection{The proof-complexity target}

By the Cook--Reckhow theorem \cite{CookReckhow79}, a propositional proof
system with polynomially bounded proof size exists if and only if
$\mathrm{NP}=\mathrm{coNP}$; hence any super-polynomial lower bound for a
sufficiently strong propositional proof system implies
$\mathrm{NP}\neq\mathrm{coNP}$ and therefore $\mathrm{P}\neq\mathrm{NP}$.
For resolution the pigeonhole principle has exponential lower bounds
\cite{Haken85}, and so it has for constant-depth Frege \cite{Ajtai94}.
Adding $\mathrm{MOD}_p$ gates defeats the Razborov--Smolensky circuit
technology \cite{Smolensky87}, and no non-trivial lower bound is known for
$\AC[2]$-Frege on any formula sequence.  In 2026 the adjacent frontier
moved: an exponential lower bound was proved, with a Lean~4 formalization,
for the bit pigeonhole principle in unrestricted dag-like resolution over
parities (Res$(\oplus)$) \cite{resxor26}, the system immediately below
$\AC[2]$-Frege, and super-linear lower bounds for random 3-CNFs became
available in bounded-depth Frege \cite{randomcnf26}.  The $\AC[2]$-Frege
lower bound for $\PHPn$ itself is the frontier open problem, and the
pseudo-solution program of \cite{Kra26} is currently its most concrete
reduction: it replaces the qualitative lower-bound question by a
quantitative probabilistic task about explicitly given objects
(Theorem~\ref{thm:chi}).  Proving anything about that task requires, first,
an exact understanding of what an adaptive decision tree with low-degree
queries can extract from the pipeline $\omega=(\rho,L)$ at the
Boolean-relevant characteristic $p=2$.  This paper supplies that
understanding for the single-variable, non-adaptive, and single-pass
product query classes; proves the answer-channel law at kernel level
(one design space, $V(n,d)^{\rho}=V(2d,d)$, with parity-locked full
rows and the i.i.d.\ law resting on the zeroing completion);
settles the unbounded-budget certification
problem at degree $1$ in the coin model exactly (Theorem~\ref{thm:floor}:
the optimal error is $2^{-\Theta(d^2)}$, attained by a non-adaptive
counting tree; under the pipeline's parity-locking the order survives
and the constant is open); proves the adjacency-certificate lemma and
reports the confirmation-chain measurements that beat the per-query
caps on the literal pipeline (Section~\ref{sec:adjacency});
restates the two-phase certification tree under the literal query
semantics, after retracting its earlier exact law as a simulator
artifact, and shows it dead against the pipeline;
reconciles the unbounded results with the printed budget of
Theorem~\ref{thm:chi}; and isolates the budgeted certification floor
over the certificate trichotomy as
the precise remaining open problem.

\subsection{Contributions}

All results concern the pipeline $\Omega(n,d)$ of \cite{Kra26} at $p=2$
with $n_\rho=2d$ (Definition~\ref{def:pipeline}).

\begin{enumerate}
\item \textbf{The answer channel law, restated at kernel level}
(Theorem~\ref{thm:channel}): matched pairs answer the constant $1$,
killed--unmatched pairs the constant $0$, free pairs answer fair bits
marginally and jointly on every \emph{row-incomplete} query set, a
fully queried free row is \emph{parity-locked} (Theorem~\ref{thm:channel}(iv)),
and the row sums answer the determined $0$ on every pigeon
(Proposition~\ref{prop:rigidity}), so aggregate row-sum queries carry
no signal at all.  The joint law rests on the zeroing completion
(Lemma~\ref{lem:completion}); the disjoint-support route recorded in an
earlier draft fails at $d\ge2$ (Remark~\ref{rem:proofroute}).  A
kernel-level verification (Remark~\ref{rem:onekernel}) proves
$V(n,d)^{\rho}=V(2d,d)$ exactly: the outer/canonical design-space
ambiguity of the previous draft dissolves, one design space carries the
printed vanishing condition, and no reading of the printed definition
realizes the coin channel beyond its row-incomplete single-variable law
(Corollary~\ref{cor:coin}).
\item \textbf{Completion of partial patterns}
(Lemma~\ref{lem:completion}): every partial assignment with fewer than
$2d$ pinned coordinates in each row extends to a design, and the number of
extensions is independent of the assigned values; a fully pinned row must
carry parity $1$.  Consequently the answer channel of any set of
single-variable queries coincides, in joint law, with an i.i.d.-coin
channel except on the event that some row's free pairs are queried in
their entirety (Corollary~\ref{cor:coin}); on that event the row is
parity-locked (Theorem~\ref{thm:channel}(iv)).
\item \textbf{Depletion}
(Lemma~\ref{lem:depletion}): conditioned on any transcript, the freeness
probability of any pair whose own cell was not probed exceeds the base
rate by a factor $1+O(s/n)$ at most.  This is the hypergeometric
negative-association bound, placed on record with the standard citations
\cite{JoagDevProschan83,DubhashiRanjan98}.
\item \textbf{Caps}: Proposition~\ref{prop:fixed} (fixed-label trees
succeed at the base rate $f$ exactly); Proposition~\ref{prop:nonadaptsingle}
(non-adaptive single-variable trees with answer-seeking labels succeed at
most $f(1+s/2)+o(1)$); Proposition~\ref{prop:nonadaptall} (non-adaptive
trees whose queries are single variables or monomials with at most
$\gamma$ distinct variable pairs each succeed at most
$(\gamma s+1)f+o(1)$ at budget $s$, with the transition
$d=\Theta_\gamma(n^{2/3}(\log k)^{-1/3})$; the earlier all-queries
generality of this cap is \emph{retracted}, with a counterexample, in
Remark~\ref{rem:retraction}); Theorem~\ref{thm:B} (all adaptive
single-variable trees of depth $e$ succeed at most
$\max(q,f)(1+O(e/n))+o(1)$, where $q=\frac{2d^2+d}{2d^2-d+n}$ is the
exact per-hit Bayes posterior; in particular the success is at most
$\frac12+o(1)$ for $2d\le\sqrt{2n}$; the cap is a \emph{budgeted}
statement, Remark~\ref{rem:Bscope}).  Hence the chi-task hypothesis of
Theorem~\ref{thm:chi} holds for these classes with margin.
\item \textbf{The AND-query no-lift}
(Theorem~\ref{thm:and}): a degree-$2$ product query is a joint test with
no per-hit advantage: its labeling posterior equals $q$ exactly, by the
polynomial identity $(f+2m)^2=f^2+4fm+4m^2$.
\item \textbf{The exact certification floor at degree $1$}
(Theorem~\ref{thm:floor}): over \emph{all} adaptive degree-$1$ trees
with no query budget, the optimal error in the coin model of
Definition~\ref{def:table} is exactly
$\mathrm{err}^*=(1-2d/n)(2d+1)(2d)!/2^{(2d+1)2d}=2^{-\Theta(d^2)}$,
attained by a non-adaptive full-scan counting tree.  The proof:
linearity of degree-$1$ answers (Lemma~\ref{lem:linear}), full-scan
dominance reducing adaptivity to the Bayes problem on the answer table
(Corollary~\ref{cor:fullscan}), sound line-count certificates
(Lemma~\ref{lem:certificates}), pigeonhole rescue
(Lemma~\ref{lem:rescue}), and the near-permutation class as the unique
obstruction, with pair posterior flat at $2d/n$
(Lemma~\ref{lem:classb}).  Verification: an exhaustive Bayes audit at
$(4,1)$ and exact enumeration of all $2^{20}$ coin matrices at $(5,2)$
(Remark~\ref{rem:floorverify}).  The rows-and-columns parity tree
(Proposition~\ref{prop:rcparity}) dominates the two-phase tree at
equal $O(n)$ cost.  The constant is a coin-model number: under the
pipeline's parity-locking the hard class vanishes and the
unbounded floor on the literal pipeline is exactly $0$
(Remark~\ref{rem:parityfloor}); the open remnant is the budgeted
all-degrees floor.
\item \textbf{Adjacency certificates and the confirmation chain}
(Section~\ref{sec:adjacency}): two adjacent $1$s certify a free pair
\emph{with certainty} (Lemma~\ref{lem:adjacency}; a matched pair
forces its row-neighbors to the determined $0$), sound against the
literal pipeline directly, with no design property used.  The
scan-then-confirm chain measures success $0.9317$ at budget $1000$ at
both $(32,2)$ and $(64,4)$ ($0.9417$ on the exact pipeline at
$(32,2)$), posterior $1.0000$ on every certified output, certificate
rate $1-e^{-B/400}$: the per-query caps are sharp, but they are
per-query bounds, and the chain is the strongest known adaptive
attacker on the pipeline, with the scaling of its constant the
sharpest open problem (Open Problem~\ref{op:adjacency}).
\item \textbf{The two-phase certification tree, corrected}
(Theorem~\ref{thm:twophase}): under the literal query semantics the
tree's error is exactly
$2^{-(2d+1)}+(1-2^{-(2d+1)})\,2^{-2d+1}$ (the count-$0$ phase-$2$
branch); the earlier exact law $1-(1-f)2^{-(2d+1)}$ is retracted as a
simulator artifact (Remark~\ref{rem:artifact}), the recorded
$0.97$-level measurements are requalified as hybrid-artifact
measurements, and the literal tree measures $0.8500$ at $(32,2)$
(Table~\ref{tab:twophase}); against the literal pipeline the tree is
dead outright (Corollary~\ref{cor:closed}).
\item \textbf{Budget reconciliation} (Section~\ref{sec:budget}):
Theorem~6.1 of \cite{Kra26} quantifies over $(d,d\log k)$-trees; every
constructed attacker is over budget, counting certificates are inert
within budget, and the open problem is the budgeted certification
floor, now organized by the certificate trichotomy (parity: dead;
counting: inert within budget; adjacency: alive,
Remark~\ref{rem:trichotomy}).
\item \textbf{Open problems} (Section~\ref{sec:open}): the
chi-hypothesis in its printed budgeted form, the budgeted
certification floor over the certificate trichotomy, the
adjacency-budget race, chaining beyond the per-hit posterior,
and the odd-characteristic analogue.
\end{enumerate}

All numerical claims are exact channel simulations or exact
enumerations over the finite coin model of Section~\ref{sec:floor}
(they sample the proved channel law or enumerate its answer tables,
not the design space), except where stated: the exact-pipeline rows of
Section~\ref{sec:adjacency} sample uniform designs over $\Des(4,2)$
directly.  The two-phase table separates two laws and
labels every row: \emph{hybrid-artifact} rows verify the artifact
simulator's certification law (Remark~\ref{rem:artifact}); \emph{literal}
rows verify Theorem~\ref{thm:twophase}.  Against the literal pipeline
the tree is dead outright: the row-sum channel answers the determined
$0$ on every pigeon (Corollary~\ref{cor:closed}), a fact now proved at
kernel level (Remark~\ref{rem:onekernel}).  The floor verification
combines an exhaustive Bayes audit at $(4,1)$, exact enumeration of
all $2^{20}$ coin matrices at $(5,2)$, and Monte Carlo at three points
in one registered run.  Measurements confirm the proved laws and are
never used in place of proofs.

\subsection{Attribution}

The pipeline, the solution condition, the reduction chain, and the
chi-task are due to J.~Kraj\'i\v{c}ek \cite{Kra26}, with the algebraic
foundation in \cite{BIKPPS97}; we quote the printed definitions of
\cite{Kra26} in Section~\ref{sec:prelim}.  Everything else in this paper
--- the answer channel law, the degree-one span lemma, the completion
lemma, the rigidity proposition, the factorization and depletion lemmas,
Propositions~\ref{prop:fixed}, \ref{prop:nonadaptsingle},
\ref{prop:nonadaptall} (in its surviving class form), Theorems~\ref{thm:B},
\ref{thm:and}, \ref{thm:twophase} (in its corrected literal form), the
two-phase tree itself, the floor lemmas and
Theorem~\ref{thm:floor}, the rows-and-columns parity tree
(Proposition~\ref{prop:rcparity}), and the adjacency certificate
(Lemma~\ref{lem:adjacency}) --- is original to this
research session and unpublished.  This revision retracts two earlier
claims of the session: the exact two-phase law $1-(1-f)2^{-(2d+1)}$
with its optimality conjecture (a simulator artifact,
Remark~\ref{rem:artifact}) and the all-degrees non-adaptive cap
(Remark~\ref{rem:retraction}).  It also supersedes the previous
draft's design-space readings discussion: the kernel-level
verification (Remark~\ref{rem:onekernel}) proves
$V(n,d)^{\rho}=V(2d,d)$ exactly, so the outer/canonical debate
dissolves, and the canonical-reading ``reversed certificate'' scenario
of the previous draft does not exist under the zeroing restriction
semantics (Remark~\ref{rem:proofroute},
Proposition~\ref{prop:rigidity}).  We work from the definitions as
printed; if the intended notion of the design set or of the tree class in
\cite{Kra26} differs from our reading, our results apply to the printed
formulation.

\subsection{Notation}\label{sec:notation}

\begin{center}
\small
\begin{tabular}{ll}
\hline
Symbol & Meaning \\
\hline
$n$ & outer dimension: $n+1$ pigeons, $n$ holes; the outer system is
$\PHPn$ \\
$d$ & degree bound of the pipeline and of the queries; $d\ge 2$ is assumed
where stated \\
$p$ & field characteristic; throughout, $p=2$ and the field is $\F_2$ \\
$x_{ij}$ & outer Boolean variable ``pigeon $i$ goes to hole $j$'',
$i\in[n+1]$, $j\in[n]$ \\
$\PHPn$ & system: hole collisions $x_{i_1j}x_{i_2j}$, two per pigeon
$x_{ij_1}x_{ij_2}$, \\
 & row sums $Q_i=1+\sum_j x_{ij}$, Boolean axioms $x_{ij}^2+x_{ij}$ \\
$\rho$ & uniform random partial injection leaving exactly $2d$ holes
unfilled \\
$D,\ R$ & the $2d+1$ free pigeons and the $2d$ free holes under $\rho$;
$F=D\times R$ the free pairs \\
$L$ & uniform random degree-$d$ design of the restricted system
(Def.~\ref{def:pipeline}) \\
$V(N,d)$ & span of the degree $\le d$ consequences of a system on $N+1$
pigeons, $N$ holes \\
$\Des(N,d)$ & affine space of linear functionals with $L(1)=1$ and
$L|_{V(N,d)}=0$ \\
$\omega=(\rho,L)$ & pipeline sample; the answer to a query $g$ is
$L(g^\rho)$ \\
$g^\rho$ & restriction of $g$ by $\rho$ (killed variables become $0$ or
the constant $1$) \\
$f$ & free-pair density $=\Prob[(i,j)\in F]=\frac{(2d+1)\,2d}{(n+1)\,n}$ \\
$m$ & matched-pair density $=\Prob[\rho(i)=j]=\frac{n-2d}{(n+1)\,n}$ \\
$q$ & per-hit Bayes posterior $=\frac{f/2}{f/2+m}=\frac{2d^2+d}{2d^2-d+n}$ \\
$s,\ e$ & number of queries (budget); depth of a tree \\
$e'$ & printed query budget of Theorem~\ref{thm:chi}:
$O(\log n)+O(d\log n)$, written $d\log k$ \\
$M$ & answer table $M(i,j)=\mathrm{ans}(x_{ij})$
(Section~\ref{sec:floor}) \\
$\mathrm{err}^*$ & the exact degree-$1$ certification floor
(Theorem~\ref{thm:floor}) \\
$\lab(\ell)$ & the pair output at leaf $\ell$ of a tree \\
$\chi(P,\rho)$ & chi-indicator of \cite{Kra26}: $1$ iff $\lab(P)\notin F$
and $1\in\Sp(W(P)^\rho)$ \\
$k,\ l$ & refutation size and $\AC$ depth in the target application;
$d=(\log k)^{O(l)}$ \\
\hline
\end{tabular}
\end{center}

All probabilities are over the uniform pipeline sample $\omega=(\rho,L)$:
$\rho$ is uniform over partial injections leaving exactly $2d$ holes free
(the free hole set $R$, the free pigeon set $D$, and the bijection
$[n+1]\setminus D\to[n]\setminus R$ being chosen independently and
uniformly), and $L$ is uniform over $\Des(2d,d)$, independently of
$\rho$.  Constants in $O(\cdot)$ and $o(\cdot)$ are absolute unless a
parameter is displayed.

\section{Preliminaries}\label{sec:prelim}

\subsection{The negated pigeonhole principle as a polynomial system}

Work over $\F_2$.  The outer system $\PHPn$ on the variables $x_{ij}$,
$i\in[n+1]$, $j\in[n]$, consists of the hole-collision polynomials
$x_{i_1j}x_{i_2j}$ ($i_1<i_2$), the two-per-pigeon polynomials
$x_{ij_1}x_{ij_2}$ ($j_1<j_2$), the row sums
$Q_i=1+\sum_{j=1}^n x_{ij}$ (in characteristic $2$, minus is plus), and
the Boolean axioms $x_{ij}^2+x_{ij}$.  The system expresses that the
assignment is a total injection of $n+1$ pigeons into $n$ holes; it is
inconsistent.  For a rectangle of $N+1$ pigeons and $N$ holes let
\[
V(N,d)\;=\;\Sp\bigl\{\,g\cdot h \;:\; h \text{ a generator},\ \deg(gh)\le
d\,\bigr\}\;\subseteq\;\F_2^{\le d}[x_{ij}]
\]
be the span of its degree-$\le d$ consequences, where $\F_2^{\le d}$
carries the multiset monomials ($x^2$ is a distinct basis element from
$x$).  A \emph{design} is a linear functional $L$ on $\F_2^{\le d}$ with
$L(1)=1$ and $L|_{V(N,d)}=0$; designs exist precisely when
$1\notin V(N,d)$, which holds in the regime of interest (Razborov's
theorem as quoted in \cite{Kra26}; verified by exact $\F_2$ elimination
in the cases of Table~\ref{tab:dims}).

\begin{table}[t]
\centering
\begin{tabular}{rrrrr}
\hline
$N$ & $d$ & $\dim \F_2^{\le d}$ & $\mathrm{rank}\,V(N,d)$ &
$\dim\Des(N,d)$\\
\hline
4 & 2 & 231 & 165 & 65\\
5 & 2 & 496 & 306 & 189\\
6 & 2 & 946 & 511 & 434\\
7 & 2 & 1653 & 792 & 860\\
8 & 2 & 2701 & 1161 & 1539\\
6 & 3 & 14190 & 12110 & 2079\\
\hline
\end{tabular}
\caption{Exact $\F_2$ dimensions of the consequence space $V(N,d)$ inside
the multiset-monomial space, and of the affine design space
$\Des(N,d)$, whose cardinality is $2^{\dim\Des}$.  In every case
$1\notin V(N,d)$ and the degree-$\le d$ closure of $V$ was verified by
exact elimination.}
\label{tab:dims}
\end{table}

\subsection{The systems ladder and where $\AC[2]$-Frege sits}

Resolution has exponential lower bounds for the pigeonhole principle
\cite{Haken85}; so do Res$(k)$ and $\AC^0$-Frege \cite{Ajtai94}.  The
system $\AC^0[p]$-Frege (constant-depth Frege with $\mathrm{MOD}_p$
gates) is equivalent to the Extended Nullstellensatz algebraic proof
system \cite{Kra19}, and the circuit lower-bound technology for
$\AC^0[p]$ \cite{Smolensky87} has never lifted to it.  In 2026 the
adjacent frontier moved: unrestricted dag-like Res$(\oplus)$ was shown
exponential for the bit pigeonhole principle \cite{resxor26}, and random
3-CNFs got super-linear bounded-depth Frege bounds \cite{randomcnf26}.
The $\AC[2]$-Frege lower bound for $\PHPn$ remains open, with no
non-trivial bound known on any formula sequence; the pseudo-solution
program of \cite{Kra26} is its most concrete current reduction.

\subsection{The reduction chain of Kraj\'i\v{c}ek}

We quote the chain as printed in \cite{Kra26}.

\begin{itemize}
\item \textbf{Theorem 2.2 of \cite{Kra26}} (from \cite{BIKPPS97}): a
$k$-step $\mathrm{F}_l(\mathrm{MOD}_p)$ refutation of $\PHPn$ yields an
Extended Nullstellensatz refutation with $S=k^{O(1)}$ extension
polynomials in $l+O(1)$ levels, accuracy $h=O(\log S)$, and degree
$d\le(\log k)^{O(l)}$.
\item \textbf{Theorems 3.2 and 3.3 of \cite{Kra26}}: if $e^{h/p}\ge 2S^2$
there is no $(d,h+\log S,S^{-1})$-solution; contrapositively, a
$((\log k)^{O(l)},\,O(\log k),\,k^{-O(1)})$-solution of $\PHPn$ rules out
$k$-step refutations.
\item \textbf{Definition 3.1 of \cite{Kra26}} (the solution condition),
quoted in Definition~\ref{def:solution} below: every tree of the
permitted class must fail, i.e.\ output a non-conflict, with probability
at least $\gamma=k^{-O(1)}$.
\item \textbf{Definition 4.3 of \cite{Kra26}} (the candidate
$\Omega(n,d)$), quoted in Definition~\ref{def:pipeline}.
\item \textbf{Theorem 6.1 of \cite{Kra26}} (the chi-task), quoted in
Theorem~\ref{thm:chi}.
\end{itemize}

\begin{definition}[pipeline $\Omega(n,d)$; Definition 4.3 of
\cite{Kra26}]\label{def:pipeline}
Fix the characteristic $p$ (here $p=2$), outer dimension $n$, and degree
$d$ with $2d\le n$.  The pipeline $\Omega(n,d)$ is the distribution on
pairs $\omega=(\rho,L)$:
\begin{enumerate}
\item $\rho$ is a uniform random partial injection from $[n+1]$ to $[n]$
leaving exactly $n_\rho=2d$ holes unfilled; let $D\subseteq[n+1]$ be the
$2d+1$ unassigned pigeons and $R\subseteq[n]$ the $2d$ unfilled holes.
The restricted system is $\neg\mathrm{PHP}$ on the rectangle $D\times R$,
still contradictory ($2d+1$ pigeons, $2d$ holes).
\item $L$ is a uniform random degree-$d$ design of the restricted
system: a linear functional on the degree-$\le d$ multiset monomials of
the variables $\{x_{ij}:(i,j)\in D\times R\}$, vanishing on the
restricted consequence space $V(n,d)^{\rho}$ (span of the restrictions
of the outer degree-$\le d$ consequences), with
\[
L(1)=1,\qquad L|_{V(n,d)^{\rho}}=0 ,
\]
an element of a non-empty affine space (designs exist by the same
separation as in Table~\ref{tab:dims}; by Remark~\ref{rem:onekernel}
the constraint space equals the canonical $V(2d,d)$, whose
complementary dimensions the table records).
\item the answer to a query $g$, a polynomial of degree $\le d$ in the
outer variables, is $\omega(g):=L(g^\rho)$, where $g^\rho$ replaces each
outer variable $x_{ij}$ by $1$ if $\rho(i)=j$, by $0$ if $i$ is assigned
and $j\ne\rho(i)$, and leaves it alone if $(i,j)\in D\times R$.  (The
printed clause leaves the cell \emph{free pigeon, matched hole} ---
$(i,j)$ with $i\in D$, $j\notin R$ --- uncovered; the \emph{zeroing
completion} sets such cells to $0$, the only choice consistent with the
determined killed--unmatched answer, Remark~\ref{rem:proofroute}.)
\end{enumerate}
A pair $(i,j)$ is \emph{free} if $(i,j)\in F:=D\times R$, \emph{matched}
if $\rho(i)=j$, and \emph{killed--unmatched} otherwise.

\emph{One design space.}  An earlier draft discussed two ``readings''
of the design set --- vanishing on $V(n,d)^{\rho}$ as printed (the
\emph{outer reading}) versus membership in the canonical $\Des(2d,d)$
of the fresh system $\neg\mathrm{PHP}$ on $D\times R$ --- and treated
their difference as material.  A kernel-level verification
(Remark~\ref{rem:onekernel}) settles the question:
$V(n,d)^{\rho}=V(2d,d)$ \emph{exactly}, so the two readings coincide
and the debate dissolves: there is only one design space, and the
operative content of the printed vanishing condition is that it forces
the free-row parities (Proposition~\ref{prop:rigidity}) while leaving
the free pairs fair on row-incomplete query sets
(Theorem~\ref{thm:channel}(iii)).  All results below are stated for
this single space.
\end{definition}

\begin{definition}[solution condition; Definition 3.1 of \cite{Kra26}]
\label{def:solution}
For parameters $(d,e,\gamma)$, a sample $\omega=(\rho,L)$ \emph{satisfies
the solution condition} if every $(d,e)$-tree $T$ --- an adaptive
decision tree of depth $e$ whose queries are polynomials of degree
$\le d$ and whose leaves are labelled by pairs --- fails on $\omega$ with
probability at least $\gamma$, where $T$ \emph{succeeds} on $\omega$ if
$\lab(T)\notin D\times R$ and $1\in\Sp(W(T)^\rho)$, with $W(T)$ the span
of the restricted system together with the query polynomials along the
path of $T$.  The family $\Omega(n,d)$ is a $(d,e,\gamma)$-pseudo-solution
if the condition holds with probability $1$.
\end{definition}

\begin{theorem}[chi-task; Theorem 6.1 of \cite{Kra26}]\label{thm:chi}
If every $(d,\,d\log k)$-tree $T'$ satisfies
\[
\Prob_{\rho,\,P}\bigl[\chi(P,\rho)=1\bigr]\;\ge\;k^{-O(1)},
\]
where $P$ is a random path of $T'$ and $\chi(P,\rho)=1$ iff
$\lab(P)\notin D^\rho\times R^\rho$ and $1\in\Sp(W(P)^\rho)$, then every
$\AC^0[l](\mathrm{MOD}_2)$-Frege refutation of $\PHPn$ has at least $k$
steps.
\end{theorem}

The conjecture of \cite{Kra26} is that the hypothesis holds even with a
constant lower bound at $k(n)=2^{n^\delta}$ for small $\delta>0$.  The
program can fail in two directions: if some tree succeeds with
probability $>1-k^{-O(1)}$, the candidate $\Omega(n,d)$ is refuted as a
pseudo-solution; if every tree's chi-probability is below $k^{-O(1)}$,
Theorem~\ref{thm:chi} cannot be applied.  The results of this paper bound
the first quantity from above for large tree classes, thereby protecting
the pseudo-solution, and isolate the exact residual open class.

\section{The answer channel at characteristic 2}\label{sec:channel}

Throughout this section $(\rho,L)\sim\Omega(n,d)$ with $p=2$ and
$2d\le n$; $d\ge2$ is assumed where stated (the channel law
(Theorem~\ref{thm:channel}) holds for $d\ge1$, the completion lemma
needs $2d\ge4$).

\subsection{The degree-one span lemma}

Recall the generators of the restricted system on $D\times R$: the
collision polynomials $x_{i_1j}x_{i_2j}$, the two-per-pigeon polynomials
$x_{ij_1}x_{ij_2}$, the row sums $Q^{c}_{i}=1+\sum_{j\in R}x_{ij}$, the
Boolean axioms $x_{ij}^2+x_{ij}$, and all their shifts by monomials of
degree at most $d-\deg$.  Write
$\mathcal{G}$ for the set of their coefficient vectors in
$\F_2^{\le d}$, so that $V(2d,d)=\Sp(\mathcal{G})$, and recall the
affine design space $\Des(2d,d)=\{L:\langle L,g\rangle=0\ \forall
g\in\mathcal{G},\ \langle L,e_0\rangle=1\}$, where $e_0$ is the constant
monomial and $\langle\cdot,\cdot\rangle$ is the coefficient pairing.

For a placement point $z_{rv}$ ($r\in D$, $v\in R$) let
$\mathrm{eval}_{z_{rv}}$ be the linear functional ``evaluate at the
partial injection placing pigeon $r$ on hole $v$ and no other pigeon'',
with $x^2$ evaluated as $x$.

\begin{lemma}[degree-one span]\label{lem:degonespan}
Assume $d\ge2$ (so each row has $2d\ge4$ cells).  The degree-$\le1$ part
of $\Sp(\mathcal{G}\cup\{e_0\})$ is exactly
$\Sp\{e_0,\;Q^{c}_1,\dots,Q^{c}_{2d+1}\}$: that is, a polynomial of
degree $\le1$ is a consequence of the restricted system (up to the
constant) if and only if it is constant on each row.
\end{lemma}

\begin{proof}
The inclusion $\supseteq$ is trivial: each $Q^c_r$ is a generator and
$e_0$ is present.  For the inclusion $\subseteq$, let $\delta$ be a
degree-$\le1$ element of $\Sp(\mathcal{G}\cup\{e_0\})$ that is not
constant on each row; we build a linear functional killing
$\mathcal{G}\cup\{e_0\}$ but not $\delta$, contradicting membership.
Choose a row $r^\ast$ on which $\delta$ is neither identically $0$ nor
identically $1$, and let $t$ be the number of cells of row $r^\ast$ on
which $\delta$ has coefficient $1$; then $1\le t\le 2d-1$.  Define a
weight vector $w$ on the cells of row $r^\ast$:
\begin{itemize}
\item if $t$ is odd, let $w$ be the all-ones vector of the row (weight
$2d$, even);
\item if $t$ is even, let $w$ be the indicator of one $1$-cell of
$\delta$ and one $0$-cell of $\delta$ (both exist since
$1\le t\le 2d-1$ and $2d-t\ge1$).
\end{itemize}
In both cases $w$ has even total weight, and the correlation
$w\cdot\delta|_{r^\ast}$ equals $t\bmod 2=1$.  Set
\[
L^{*}\;=\;\sum_{v\in R} w_{v}\,\mathrm{eval}_{z_{r^\ast v}} .
\]
Every generator vanishes pointwise at every placement point (collisions
and two-per-pigeon need two placed pigeons on a common hole or row;
Booleans vanish on $0$/$1$ points), hence their shifts do as well, so
$\langle L^{*},g\rangle=0$ for every shifted collision, two-per-pigeon,
and Boolean generator.  For the row sums,
$Q^{c}_u(z_{r^\ast v})=[u\neq r^\ast]$, whence
$\langle L^{*},Q^c_u\rangle=[u\neq r^\ast]\cdot\sum_v w_v
=[u\neq r^\ast]\cdot 0=0$, the total weight of $w$ being even; likewise
$\langle L^{*},e_0\rangle=0$.  Thus $L^{*}$ annihilates
$\mathcal{G}\cup\{e_0\}$, while
$\langle L^{*},\delta\rangle=\sum_v w_v\delta_{r^\ast v}=1$.  A
polynomial annihilated by no such functional cannot lie in the span;
contradiction.  Finally, the polynomials constant on each row are exactly
the combinations $c\,e_0+\sum_r a_r Q^c_r$ (row-$r$ coefficient
$c+a_r$), which proves the characterization.

For the pipeline's design space the constraint set is the span of the
restrictions of the outer generators; by Remark~\ref{rem:onekernel}
($V(n,d)^{\rho}=V(2d,d)$ exactly) that span coincides with
$\Sp(\mathcal{G})$, so the proof above covers the pipeline directly:
the degree-$\le1$ part of the constraint span (with $e_0$) is
$\Sp\{e_0,\;Q^{c}_1,\dots,Q^{c}_{2d+1}\}$, and the row parities are
forced to $1$ on the free rows, the assigned rows carrying no free
coordinates at all (their variables are zeroed by the restriction).
\end{proof}

\subsection{The channel law}

\begin{theorem}[answer channel law, single variables; the kernel
law]\label{thm:channel}
Let $(i,j)$ be an outer pair and $a(x_{ij})=L(x_{ij}^\rho)$ the answer to
the single-variable query $x_{ij}$.
\begin{enumerate}
\item[(i)] If $\rho(i)=j$ (matched), then $a=1$ deterministically.
\item[(ii)] If $i\notin D$ and $j\ne\rho(i)$ (killed--unmatched), then
$a=0$ deterministically.
\item[(iii)] If $(i,j)\in F$ (free), then $a$ is uniform on $\{0,1\}$;
moreover, on any \emph{row-incomplete} query set $S$ --- no row of the
rectangle has all $2d$ of its cells in $S$ --- the answers of the free
pairs of $S$ are jointly uniform: every pattern occurs for the same
number of designs (Lemma~\ref{lem:completion}(a)).
\item[(iv)] \emph{Parity-locking.}  If the $2d$ cells of a free row are
all queried, the answer pattern is locked to odd parity: exactly
$2^{2d-1}$ of the $2^{2d}$ patterns occur, each for the same number of
designs (Lemma~\ref{lem:completion}(b)); the row XOR is $1$ for every
occurring pattern, and the all-zero pattern never occurs.
\end{enumerate}
\end{theorem}

\begin{proof}
(i) If $\rho(i)=j$ then $x_{ij}^\rho=1$ and $a=L(1)=1$ by the design
axiom.  (ii) If $i\notin D$ and $j\ne\rho(i)$ then $x_{ij}^\rho=0$ and
$a=0$.  (iii) Marginal uniformity: the coordinate $L\mapsto L(x_{ij})$ is
constant on the affine design space if and only if the vector
$e_{ij}$ lies in the orthogonal complement of its direction space,
which is $\Sp(\mathcal{G}\cup\{e_0\})$ (the restricted generators
together with the constant).  By
Lemma~\ref{lem:degonespan} every degree-$\le1$ element of that span is
constant on each row; $e_{ij}$ is supported on the single cell $(i,j)$
and its row has $2d\ge2$ cells, so it is not row-constant and not the
constant vector.  Hence the coordinate is non-constant, and a
non-constant linear coordinate on an affine space over $\F_2$ takes the
values $0$ and $1$ equally often: $a$ is a fair bit.  The joint
statement is Lemma~\ref{lem:completion}(a): the zeroing completion
exhibits, for every row-incomplete pattern $\psi$ on $S$, exactly
$2^{\dim\Des(2d,d)-|S|}$ designs, independently of $\psi$, so the
answers of the free pairs of $S$ are jointly uniform.  (iv) If $S$
contains a full free row $r$, Lemma~\ref{lem:completion}(b) applies:
extensions exist exactly for the pinned patterns of parity $1$ on row
$r$ (rigidity \eqref{eq:rowparity}), and every admissible pattern
occurs for the same number of designs, so the answer law is uniform on
the $2^{2d-1}$ odd-parity patterns and gives nothing to the rest: the
row is parity-locked, its XOR equal to $1$.
\end{proof}

\begin{remark}[proof-route correction: no disjoint supports at
$d\ge2$]\label{rem:proofroute}
An earlier recorded route to the joint law (iii) argued from a kernel
basis with pairwise-disjoint single-variable supports.  That route
fails at $d\ge2$ and is superseded: the kernel's single-variable
projection is the even-parity subspace of each free row, and within a
row of $2d$ coordinates no $2d-1$ even-weight vectors with pairwise
disjoint supports exist, their total weight being at least
$2(2d-1)>2d$ (the counting obstruction; at the restricted point
$(4,2)$ the reduced basis already has $15$ overlapping single-variable
pairs).  At $(2,1)$ the disjoint route happens to work; at every $d$
the correct route is the zeroing completion, Lemma~\ref{lem:completion},
on which (iii) and (iv) rest.  The completion itself is forced: the
printed definition of $g^\rho$ leaves the cell (free pigeon, matched
hole) uncovered, and any completion leaving those variables symbolic
would contradict the determined killed--unmatched answer (ii).
\end{remark}

\begin{proposition}[row-parity rigidity; the aggregate channel is
closed]\label{prop:rigidity}
Every design $L$ of the restricted system satisfies
\begin{equation}\label{eq:rowparity}
\sum_{j\in R}L(x_{ij})=1\qquad\text{for every row } i\in D .
\end{equation}
Consequently, in the pipeline, the outer row-sum query
$Q_i=1+\sum_jx_{ij}$ answers the determined $0$ on \emph{every} pigeon:
for a free pigeon $i$, $Q_i^{\rho}=Q^{c}_{\mathrm{row}(i)}$ is itself a
restricted consequence; for an assigned pigeon $i$ with $\rho(i)=j_0$,
the zeroing restriction gives
$Q_i^{\rho}=1+x_{i,j_0}^{\rho}=1+1=0$ identically.  Row queries
therefore carry no information at all, and no certificate built on
$Q_i$ exists against the pipeline.  (An earlier draft derived, under a
canonical reading, a ``reversed certificate'' with assigned pigeons
answering $1$; that computation treated assigned-row variables as
symbolic, which the zeroing semantics excludes, and the scenario does
not exist --- the kernel-level verification confirms
$L(Q_i^{\rho})=0$ under the single design space,
Remark~\ref{rem:onekernel}.)
\end{proposition}

\begin{proof}
The restricted row sum $Q^c_i=1+\sum_{j\in R}x_{ij}$ is a generator of
degree $1\le d$, so $L(Q^c_i)=0$; expanding with $L(1)=1$ gives
\eqref{eq:rowparity}.  For the outer query: if $i\in D$ then
$Q_i^{\rho}=1+\sum_{j\in R}x_{ij}=Q^c_i$, whence $L(Q_i^{\rho})=0$.  If
$i\notin D$ with $\rho(i)=j_0$, the zeroing completion kills every cell
of row $i$ except the matched one, so
$Q_i^{\rho}=1+1+\sum_{j}0=0$ as a polynomial, whence
$L(Q_i^{\rho})=0$.  (The zeroing completion is the only one consistent
with the determined killed--unmatched answer;
Remark~\ref{rem:proofroute}.)
\end{proof}

\begin{remark}[the one-kernel theorem]\label{rem:onekernel}
Exact $\F_2$ elimination (every generator row built through the outer
route --- outer polynomials, outer shifts, then restriction --- and
compared with the independently built canonical list by rank and by
containment in both directions) gives $V(n,d)^{\rho}=V(2d,d)$ at every
tested point, with the ranks matching the canonical values of
Table~\ref{tab:dims}:
\begin{center}
\small
\begin{tabular}{rrrrl}
\hline
outer $(n,d)$ & restricted $(2d,d)$ & $\mathrm{rank}\,V(n,d)^{\rho}$ &
$\mathrm{rank}\,V(2d,d)$ & containment\\
\hline
(7,1)  & (2,1) & 3     & 3     & both ways\\
(15,2) & (4,2) & 165   & 165   & both ways\\
(31,3) & (6,3) & 12110 & 12110 & both ways\\
\hline
\end{tabular}
\end{center}
The reason is structural: an assigned pigeon's $Q_i$ restricts to the
zero polynomial ($1+1=0$), the killed Booleans, collisions, and
two-hole monomials restrict to $0$, and $\deg(h^{\rho})=\deg(h)$
whenever $h^{\rho}\neq0$, so every restricted outer generator is a
canonical generator shift or $0$.  Three consequences are used below.
First, the outer/canonical design-space debate of the previous draft
dissolves: there is only one kernel, and the corpus's adoption of the
outer reading is moot as a definitional matter.  Second, the printed
vanishing condition itself forces the free-row parities
(Proposition~\ref{prop:rigidity}).  Third, no design space satisfying
the printed definition realizes the coin channel beyond its
row-incomplete single-variable law: the ``outer-design kernel'' that
masks the row parities in the never-queried killed coordinates
violates the printed condition for about half its designs (measured
$50382/100000$ and $49962/100000$), so the coin channel is exactly the
row-incomplete single-variable answer law, the precise content of
Corollary~\ref{cor:coin} with its $E_{\mathrm{full}}$ clause.
\end{remark}

\subsection{Completion of partial answer patterns}

\begin{lemma}[completion]\label{lem:completion}
Assume $d\ge2$.  Let $S$ be a set of cells of the rectangle $D\times R$
and $\psi:S\to\F_2$.
\begin{enumerate}
\item[(a)] If no row contains all $2d$ of its cells in $S$, then $\psi$
extends to a design: there is $L\in\Des(2d,d)$ with $L|_S=\psi$; and the
number of designs with $L|_S=\psi$ equals $2^{\dim\Des(2d,d)-|S|}$,
independently of $\psi$.
\item[(b)] If $S$ pins all cells of some row $r$, extensions exist only
if the pinned values on row $r$ have the parity forced by rigidity
(parity $1$ for a free row; assigned rows carry no free cells whose
design values could be pinned).  If every fully pinned row meets its
parity requirement,
extensions exist, and their number is
$2^{\dim\Des(2d,d)-|S|+\{\#\text{ fully pinned rows}\}}$,
independently of the admissible $\psi$.  In the applications below only
free cells are ever pinned (killed cells answer determinedly and impose
no design constraint), so every pinned row is a free row.
\end{enumerate}
\end{lemma}

\begin{proof}
The number of designs with $L|_S=\psi$ equals
$2^{\,n-\mathrm{rank}(A)-|S|+\rho_S}$, where $n$ is the number of
monomials, $A$ the constraint system $(\mathcal{G},e_0)$ of rank
$\mathrm{rank}(A)$, and $\rho_S$ the dimension of the intersection of
$\Sp\{e_u:u\in S\}$ with $\Sp(\mathcal{G}\cup\{e_0\})$; the solution set,
when non-empty, is an affine space over the nullspace, giving the
$\psi$-independence.  So it suffices to compute $\rho_S$ and feasibility.

(a) Suppose no row is fully pinned.  If
$\sum_u\delta_ue_u\in\Sp(\mathcal{G}\cup\{e_0\})$ for some $\delta$
supported on $S$, then by Lemma~\ref{lem:degonespan} the polynomial
$\sum_u\delta_ux_u$ is constant on each row, say with value $c+a_r$ on
row $r$; but every row has a cell outside $S$, where the polynomial
vanishes, forcing $c+a_r=0$ for every $r$, i.e.\ the polynomial is
identically zero and $\delta=0$.  Hence $\rho_S=0$, the $S$-equations
are independent of $(\mathcal{G},e_0)$, every $\psi$ is feasible, and
the count is $2^{\dim\Des-|S|}$, where
$\dim\Des=n-\mathrm{rank}(A)-1$.

(b) Suppose $S$ pins the full rows of a set $R^{\ast}$ of rows, and
parity $1$ holds on each of them.  Any dependency now satisfies: on a
non-fully-pinned row the polynomial vanishes off $S$, forcing
$c+a_r=0$ there; on a fully pinned row $r\in R^{\ast}$ the polynomial
equals the constant $c+a_r$.  Hence the dependency space is spanned by
$\{\,\sum_{u\in\mathrm{row}(r)}e_u \;:\; r\in R^{\ast}\,\}$, of
dimension $|R^{\ast}|$, since $\sum_{u\in\mathrm{row}(r)}x_{u}
=Q^c_r+1\in\Sp(\mathcal{G}\cup\{e_0\})$ and no smaller combination lies
in the span (the restriction of the dependency to a single fully pinned
row is the all-ones pattern, which is not a dependence unless the whole
row is used).  Thus
$\rho_S=|R^{\ast}|$ and the count is
$2^{\dim\Des-|S|+|R^{\ast}|}$ for every feasible $\psi$.  Feasibility:
rigidity \eqref{eq:rowparity} forces the parity condition, so it is
necessary.  For sufficiency, remove one pinned cell from each fully
pinned row, obtaining $S'$ with no full row; by (a) there is a design
$L'$ with $L'|_{S'}=\psi|_{S'}$, and by \eqref{eq:rowparity} the value
of $L'$ on each removed cell is forced to equal
$1-\sum_{\text{rest of its row}}\psi$, which is exactly the prescribed
$\psi$-value by the parity hypothesis.  Hence $L'$ extends $\psi$.
\end{proof}

\subsection{The coin channel and its equivalence with the pipeline for
single-variable trees}

\begin{definition}[coin channel]\label{def:coin}
$\mathsf{Coin}(n,d)$ is the channel with the same status process $\rho$
as the pipeline, in which each free pair carries an independent fair bit,
independent across pairs and of $\rho$, a matched pair answers $1$ and a
killed--unmatched pair answers $0$, and the answer to a product
$g=\prod x_{i_tj_t}$ is the $\F_2$-product of the answers to its
factors.
\end{definition}

\begin{corollary}[coin-channel equivalence for single-variable
trees]\label{cor:coin}
Let $T$ be an adaptive tree of depth $e$ all of whose queries are single
outer variables, and let $E_{\mathrm{full}}$ be the event (under
$\Omega$ or under $\mathsf{Coin}$, identically) that the set of free
pairs queried along the path of $T$ contains all $2d$ cells of some row
of the rectangle.  Then
\[
\Bigl|\Prob_{\Omega}\bigl[T \text{ outputs a free pair}\bigr]
-\Prob_{\mathsf{Coin}}\bigl[T \text{ outputs a free pair}\bigr]\Bigr|
\;\le\;\Prob\bigl[E_{\mathrm{full}}\bigr]
\;\le\;\min\Bigl\{1,\;\frac{(n+1)\binom{e}{2d}}{\binom{n}{2d}}\Bigr\} .
\]
\end{corollary}

\begin{proof}
Conditionally on $\rho$ and on the complement of
$E_{\mathrm{full}}$, Lemma~\ref{lem:completion}(a) gives that the
probability, over the design space, of any fixed
answer pattern on the queried free pairs equals $2^{-a(\rho)}$, where
$a(\rho)$ is the number
of queried free pairs, independently of the pattern; the killed pairs
answer determinedly in both models.  This is exactly the coin-channel
likelihood.  Equal likelihoods given $\rho$ imply equal joint laws of
$(\rho,\text{transcript})$ off $E_{\mathrm{full}}$, hence equal output
laws off that event; on the event the two laws may differ, and the total
variation between the output distributions is at most the probability of
the event.  For the bound: if the path queries $s_r$ cells of row $r$,
then the row's free pairs are all queried only if $s_r\ge2d$ and the
row's free set is a $2d$-subset of the queried cells; by symmetry of
$\rho$, a fixed $2d$-set is the free set of a fixed row with probability
$\binom{n}{2d}^{-1}$, and a union bound over rows and over the at most
$\binom{e}{2d}$ subsets of the queried cells gives the claim.
\end{proof}

Thus for the single-variable class the coin channel is not an
idealization but an exact law up to the negligible event
$E_{\mathrm{full}}$; in the parameter regime of interest,
$e=d\log k\ll n$, the bound is $o(1)$ (and at $(n,d,e)=(32,2,33)$ it is
below $4\times10^{-3}$).  The kernel-level verification sharpens the
statement to its exact form: off $E_{\mathrm{full}}$ the total
variation between the two answer laws is exactly $0$ (on every
marginal and on every row-incomplete joint, confirmed at the points of
Remark~\ref{rem:onekernel}); on a fully queried free row it is exactly
$\tfrac12$, the design law being uniform on the $2^{2d-1}$ odd
patterns against the coin law's $2^{2d}$; and the coin channel is
realized by no design space satisfying the printed definition
(Remark~\ref{rem:onekernel}).  The coin-channel samplers accordingly
over-sample: a coin sample satisfies the design constraints with
probability $\approx2^{-(2d+1)}$ (measured $0.12587$, $0.03051$,
$0.00793$ against $2^{-3},2^{-5},2^{-7}$ at the restricted points
$(2,1),(4,2),(6,3)$).

\section{Fixed labels and non-adaptive trees}\label{sec:nonadaptive}

Throughout, $f=\Prob[(i,j)\in F]=\frac{(2d+1)2d}{(n+1)n}$ and
$m=\Prob[\rho(i)=j]=\frac{n-2d}{(n+1)n}$; note $f=O(d^2/n^2)$ and
$f+m\le 2/n$ whenever $2d\le\sqrt{n}$.

The following trivial but fundamental disintegration underlies all caps.
Recall that a \emph{status pattern} $\sigma$ assigns to each cell of a
set $S$ one of the labels free, matched, killed--unmatched.

\begin{lemma}[factorization]\label{lem:factor}
Let $A$ be any event determined by the answers to a fixed set of queries,
and let $\sigma(A)$ denote the $\sigma$-algebra generated by the status
patterns of the queried cells that are consistent with those answers.
For every event $B$ depending only on $\rho$ (such as $p\in F$):
\[
\Prob[B\mid A]\;=\;\sum_{\sigma}\Prob[\sigma\mid A]\,
\Prob[B\mid\sigma] .
\]
\end{lemma}

\begin{proof}
Given a status pattern $\sigma$, the answers are a deterministic function
of the status pattern and the design values of the free queried
coordinates; the latter are, conditionally on $\sigma$, independent of
$\rho$ (the design is sampled independently of $\rho$, and $\sigma$ is
$\rho$-measurable).  Hence $\Prob[B\mid\sigma,\text{answers}]
=\Prob[B\mid\sigma]$ for every $\rho$-event $B$, and the tower rule
applied to the partition generated by $(\sigma,\text{answers})$ gives
the claim.
\end{proof}

\begin{lemma}[hypergeometric depletion]\label{lem:depletion}
Let $S$ be a set of at most $s$ cells and let $\sigma$ be a status
assignment on $S$.  Then for every cell $p\notin S$,
\[
\Prob\bigl[\,p\in F \,\big|\, \mathrm{status}(S)=\sigma\,\bigr]
\;\le\;\frac{(2d+1)}{(n+1-s)}\cdot\frac{2d}{(n-s)}
\;=\;f\,\Bigl(1+O\bigl(\tfrac{s}{n}\bigr)\Bigr)
\qquad (s\le n/2),
\]
and the same bound holds conditionally on any event determined by the
answers, by Lemma~\ref{lem:factor}.  The pigeon-freeness indicators
$\{[i\in D]\}_i$ and the hole-freeness indicators $\{[j\in R]\}_j$ are
inclusion indicators of without-replacement samples and are negatively
associated in the sense of Joag-Dev and Proschan
\cite{JoagDevProschan83} (their Section~3 and Theorem~8); the pair
indicators $[p\in F]=[i\in D]\cdot[j\in R]$ inherit the negative
association by the closure of negative association under products over
disjoint supports, in the occupancy formulation of Dubhashi and Ranjan
\cite{DubhashiRanjan98}.
\end{lemma}

\begin{proof}
Reveal the statuses on $S$; let $r_p$ be the number of pigeons other
than $p$'s pigeon revealed free, and $r_h$ the number of holes revealed
free.  The completions of $\sigma$ to full partial injections form a
without-replacement scheme: the free pigeons of the completion are a
uniform subset of the right size, and similarly for holes.  Direct
counting gives
$\Prob[i\in D\mid\sigma]=\frac{2d+1-r_p}{n+1-s'}$ and
$\Prob[j\in R\mid\sigma]=\frac{2d-r_h}{n-s'}$ with $s'\le s$ the number
of cells of $S$ touching the relevant coordinate, and the joint
probability satisfies
\[
\Prob[p\in F\mid\sigma]
\;\le\;\frac{2d+1}{n+1-s}\cdot\frac{2d}{n-s}
\;=\;\frac{(2d+1)\,2d}{(n+1)n}
 \cdot\frac{(n+1)n}{(n+1-s)(n-s)}
\;=\;f\bigl(1+O(s/n)\bigr),
\]
the joint inequality being the negative dependence of sampling without
replacement: the freeness indicators of the pigeons are the inclusion
indicators of a uniform without-replacement sample, a canonical
negatively associated family \cite{JoagDevProschan83}, and the product
form for $p\in F$ follows from the closure under products of
negatively associated families over disjoint index sets
\cite{JoagDevProschan83}, in the occupancy formulation of
\cite{DubhashiRanjan98}.  The final sentence is
Lemma~\ref{lem:factor} with $B=\{p\in F\}$.
\end{proof}

\begin{proposition}[fixed labels]\label{prop:fixed}
Let $T$ be any tree, with queries of any degree and any depth, all of
whose leaves carry the same label $(i_0,j_0)$.  Then its success
probability (output a free pair on a live path) equals exactly
\[
\Prob[(i_0,j_0)\in F]\;=\;\frac{(2d+1)}{(n+1)}\cdot\frac{2d}{n}\;=\;f,
\]
independent of the queries, the depth, and the answers.
\end{proposition}

\begin{proof}
Success is the event that the fixed output pair is free, which depends
on $\rho$ alone.  The pipeline is invariant under the full pigeon and
hole permutation groups (the design space $\Des(2d,d)$ is invariant, $L$
being uniform over it), so every one of the $(n+1)n$ pairs is free with
the same probability; the expected number of free pairs is
$|D|\,|R|=(2d+1)2d$, so the common value is
$(2d+1)2d/((n+1)n)=f$.  The queries are irrelevant because the label is
fixed.
\end{proof}

\begin{proposition}[non-adaptive single-variable trees with answer-seeking
labels]\label{prop:nonadaptsingle}
Let $T$ be a non-adaptive tree with single-variable queries
$q_1,\dots,q_s$ fixed in advance and the leaf label a function of the
$s$ answers.  Assume the label is either a pair that answered $1$, or,
if all answers are $0$, an unqueried pair.  Then
\[
\Prob[T\text{ succeeds}]
\;\le\;\frac{s\,f}{2}\;+\;f\Bigl(1+O\bigl(\tfrac{s}{n}\bigr)\Bigr)+o(1)
\;=\;f\bigl(1+\tfrac{s}{2}\bigr)+O\bigl(\tfrac{sf}{n}\bigr)+o(1),
\]
and the error is at least $\frac12-o(1)$ whenever $s\le n/2$ and
$2d\le\sqrt{n/2}$.
\end{proposition}

\begin{proof}
Let $A_1$ be the event that some answer equals $1$.  On $A_1$ the label
is a pair that answered $1$; success requires some queried pair to
answer $1$ while free, and the label can coincide with at most one of
them, so by the union bound and linearity (no independence needed; each
$q_t$ is a fixed pair)
\[
\Prob[\text{success},\,A_1]
\;\le\;\sum_{t=1}^{s}\Prob\bigl[q_t\in F \text{ and } a(q_t)=1\bigr]
\;=\;\sum_{t=1}^{s}\Prob[q_t\in F]\cdot\tfrac12
\;=\;\frac{sf}{2},
\]
using Theorem~\ref{thm:channel}(iii) and symmetry.  On the complement
$A_0$ (all answers $0$) the label is an unqueried pair $p$ determined by
the all-zero pattern, and Lemma~\ref{lem:depletion} applied to the $s$
queried cells gives
$\Prob[p\in F\mid A_0]\le f(1+O(s/n))$.  Summing gives the bound.  The
numeric claim: $f(1+s/2+O(s/n))\le \frac{4d^2}{n^2}\bigl(1+\frac{n}{4}
+O(1)\bigr)\le\frac{d^2}{n}+o(1)\le\frac12$ for $2d\le\sqrt{n/2}$ and
$s\le n/2$.
\end{proof}

\begin{remark}
Without the fallback discipline (a tree may output a queried pair that
answered $0$, which is free with probability $f/2$ per query), the
$s f/2$ term doubles to $s f$; that case is covered by
Proposition~\ref{prop:nonadaptall} within the variable/monomial class.
The measured adaptive first-hit
success $0.2367$ at $(n,d,s)=(32,2,33)$ sits below the bound
$f(1+s/2)=0.256$, as it must.
\end{remark}

\begin{proposition}[non-adaptive variable/monomial-class trees]
\label{prop:nonadaptall}
Let $T$ be a non-adaptive tree with $s\le n/2$ fixed queries, each of
which is either a single variable or a monomial in at most $\gamma$
distinct variable pairs (so $\gamma=1$ for single-variable queries and
$\gamma=2$ for degree-$2$ monomials), the label an arbitrary function
of the answers, possibly an unqueried pair.  Then
\[
\Prob[T\text{ succeeds}]
\;\le\;(\gamma s+1)\,f\,\Bigl(1+O\bigl(\tfrac{s}{n}\bigr)\Bigr)+o(1).
\]
Consequently, at the program's budget $s=e'=d\log k$, the chi-task error
requirement $k^{-O(1)}$ is safe within this class as long as
$(\gamma s+1)f\le 1-k^{-O(1)}$, i.e.\ up to
$d=\Theta_\gamma\bigl(n^{2/3}(\log k)^{-1/3}\bigr)$; beyond that the
bound ceases to decide the hypothesis.
\end{proposition}

\begin{proof}
Let $Q$ be the set of variable pairs occurring in the queries,
$|Q|\le\gamma s$.  Split by whether the output is in $Q$:
\[
\Prob[\text{success}]
=\Prob[\text{success},\,\lab\in Q]
 +\Prob[\text{success},\,\lab\notin Q].
\]
Queried output: for each $q\in Q$,
$\Prob[\lab=q,\,q\in F]\le\Prob[q\in F]=f$ by symmetry
(Proposition~\ref{prop:fixed}); the union bound gives the first term at
most $|Q|f\le\gamma s f$.

Unqueried output: condition on the full answer pattern $\pi$; the label
is a fixed cell $p(\pi)\notin Q$.  The answers reveal information about
$\rho$ only through the status pattern of the queried cells
(Lemma~\ref{lem:factor}: for variables and for monomials under the
product semantics the answer of each query is a deterministic function
of the answers to its at most $\gamma$ component pairs; on the literal
pipeline the monomial answer is instead a fresh value of the same data,
and Lemma~\ref{lem:factor} needs only that each answer be a function of
the queried cells' statuses and design values, the latter independent
of $\rho$), and by
Lemma~\ref{lem:depletion} every status assignment $\sigma$ satisfies
$\Prob[p(\pi)\in F\mid\sigma]\le f(1+O(s/n))$; averaging over
$\sigma$ with the weights $\Prob[\sigma\mid\pi]$ gives
$\Prob[p(\pi)\in F\mid\pi]\le f(1+O(s/n))$, and averaging over $\pi$
gives the second term.

At $s=d\log k$, $(\gamma s+1)f=\Theta(\gamma d^3\log k/n^2)$, which is
at most $1-k^{-O(1)}$ exactly up to
$d=\Theta_\gamma(n^{2/3}(\log k)^{-1/3})$.
\end{proof}

\begin{remark}[retraction of the all-queries generality]
\label{rem:retraction}
Proposition~\ref{prop:nonadaptall} was previously claimed for
arbitrary queries of any degree, with the bound $(s+1)f+o(1)$ and the
transition $d>\Theta(n^{2/3}(\log k)^{-1/3})$ read as an all-degrees
statement.  That generality is \textbf{false} and is retracted.  The
rc counting tree --- non-adaptive, $s=2n+1$: query all $n+1$ row sums
$Q_i$ and all $n$ column sums $K_j=1+\sum_ix_{ij}$, certify every line
whose count of $1$s differs from $1$, and output a certified row's
pigeon paired with a certified column's hole --- achieves exact success
$1-3\cdot 2^{-(2d+1)}$ at $s=2n+1$ queries: $0.9067$ at $(64,2)$,
against the retracted cap's $(s+1)f=130f=0.625$ at the same point.
The retracted proof's per-query accounting implicitly required the
answer-$1$ event of a query to carry $O(f)$ of free mass; that holds
for single variables and bounded-degree monomials, whose answer-$1$
event entails all $\gamma$ component pairs to be free or matched, but
fails for linear forms: the row sum touches the $2d$ free cells of its
row and answers $1$ with probability $\approx\frac12\gg f$.  The
counting certificates that break the general claim are exactly the
line-count certificates of Lemma~\ref{lem:certificates}
(Section~\ref{sec:floor}), where they are proved sound and shown to be
optimal.
\end{remark}

The argument uses only pair symmetry, the determined-answer law, and
depletion, so Proposition~\ref{prop:nonadaptall} holds at every prime
$p$ and under either semantics for product queries, within the
variable/monomial class.

\section{Adaptive single-variable trees}\label{sec:adaptive}

\begin{lemma}[own-answer Bayes bound]\label{lem:bayes}
In the coin channel, let $p$ be a pair queried $h\ge1$ times along a
path (necessarily with the identical answer repeated, since the same
design bit is re-read).  Then
\[
\Prob\bigl[p\in F\mid a_1,\dots,a_h\bigr]\;\le\;\max\bigl(q,\,f\bigr),
\qquad
q=\frac{f/2}{f/2+m}=\frac{2d^2+d}{2d^2-d+n},
\]
with equality throughout for a single answer $1$.  The same bound holds
in the pipeline up to the event $E_{\mathrm{full}}$
(Corollary~\ref{cor:coin}).
\end{lemma}

\begin{proof}
Given only the pair's own answers, its status is matched (all answers
$1$, likelihood $1$), free (the single design bit, likelihood
$\tfrac12$ for any consistent pattern), or killed--unmatched (all
answers $0$, likelihood $1$); the priors are $m$, $f$, $1-f-m$.  Bayes
gives
\[
\Prob[p\in F\mid a_1=\dots=a_h=1]
=\frac{f/2}{f/2+m}=q,
\quad
\Prob[p\in F\mid \text{some } a_t=0]
=\frac{f/2}{f/2+1-f-m}\le f .
\]
The transfer to the pipeline is Corollary~\ref{cor:coin}.
\end{proof}

\begin{lemma}[one-cell perturbation]\label{lem:perturb}
In the coin channel, fix a cell $p$ and a partial answer sequence $T'$
for the other queried cells ($|T'|\le s$).  Then for every pair of
statuses $s_1,s_2\in\{F,M,U\}$ of $p$,
\[
\frac{\Prob[T'\mid \sigma_p=s_1]}{\Prob[T'\mid \sigma_p=s_2]}
\;\in\;\bigl[\,1-C\tfrac{s}{n},\;1+C\tfrac{s}{n}\,\bigr]
\qquad\text{whenever } s\le n/2,
\]
for an absolute constant $C$, where $\sigma_p$ denotes the status of
$p$.  The same holds in the pipeline off the event
$E_{\mathrm{full}}$ of Corollary~\ref{cor:coin}.
\end{lemma}

\begin{proof}
Reveal $T'$ answer by answer, and condition throughout on the revealed
prefix and on $\sigma_p$.  Conditional on the prefix, the next answer's
law is determined by the conditional composition of the pools: the
probability that the next queried cell is matched, free, or
killed--unmatched is a hypergeometric quantity in the remaining pigeon
and hole pools.  The two hypotheses $\sigma_p=s_1$ and $\sigma_p=s_2$
differ by the position of one fixed pigeon--hole pair (free in one,
matched or absent in the other), and by the elementary finite-population
count of Lemma~\ref{lem:depletion} every conditional status probability
of an as-yet-unqueried cell differs between the two branches by a factor
$1+O(1/n)$ (one consumed or freed pair shifts the numerators and
denominators of the same hypergeometric ratios by at most one).
Conditional on the prefix and the status of the next cell, its answer is
either determined identically in both branches, or it is a fresh fair
design bit in both branches; hence each per-step likelihood ratio lies
in $[1-C/n,\,1+C/n]$, and after at most $s$ steps the ratio lies in
$[1-Cs/n,\,1+Cs/n]$.  The pipeline statement is
Corollary~\ref{cor:coin}: off $E_{\mathrm{full}}$ the two channels have
equal likelihoods for single-variable transcripts.
\end{proof}

\begin{theorem}[adaptive single-variable cap]\label{thm:B}
Let $T$ be an adaptive tree of depth $e$ all of whose queries are single
outer variables.  Then, in the pipeline at $p=2$,
\[
\Prob[T\text{ succeeds}]
\;\le\;\max\bigl(q,\,f\bigr)
\Bigl(1+O\bigl(\tfrac{e}{n}\bigr)\Bigr)+o(1),
\qquad
q=\frac{f/2}{f/2+m}=\frac{2d^2+d}{2d^2-d+n}
\;\approx\;\frac{2d^2}{2d^2+n}.
\]
In particular, for $2d\le\sqrt{2n}$ we have $q\le\frac12+o(1)$ and
$f=o(1)$, so every such tree fails with probability at least
$\frac12-o(1)\ge k^{-O(1)}$ for every polynomial $k$: the chi-task
hypothesis of Theorem~\ref{thm:chi} holds for the entire adaptive
single-variable class at $p=2$ in the super-polynomial regime.
\end{theorem}

\begin{proof}
By Corollary~\ref{cor:coin} it suffices to prove the bound in the coin
channel and add the $E_{\mathrm{full}}$ slack, which is absorbed into
$o(1)$.  Decompose success over the leaves $\ell$:
\[
\Prob[\text{success}]
=\sum_{\ell}\Prob[\mathrm{reach}(\ell)]\,
\Prob\bigl[\lab(\ell)\in F \mid \mathrm{reach}(\ell)\bigr],
\]
the reach events being disjoint.  Fix a leaf with transcript
$T_\ell$ (queries and answers along the path) and label $p$.

Case 1: $p$ occurs among the path's queries.  Let $a(p)$ be its own
answers and $T'$ the remaining answers.  Condition on the status
$\sigma_p$ of $p$; by Lemma~\ref{lem:factor} (the design values carry no
information about $\rho$ beyond the statuses), the posterior odds
between any two statuses of $p$ given the full transcript equal the
prior odds times the own-answer likelihood ratio times the perturbation
factor of Lemma~\ref{lem:perturb}:
\[
\frac{\Prob[\sigma_p=s_1\mid T_\ell]}{\Prob[\sigma_p=s_2\mid T_\ell]}
=\frac{\Prob[\sigma_p=s_1]}{\Prob[\sigma_p=s_2]}
\cdot\frac{\Prob[a(p)\mid s_1]}{\Prob[a(p)\mid s_2]}
\cdot\Theta ,
\qquad
\Theta\in\bigl(1-O\bigl(\tfrac{e}{n}\bigr),\,1+O\bigl(\tfrac{e}{n}\bigr)\bigr),
\]
with priors $f$, $m$, $1-f-m$ for $F$, $M$, $U$.  If $a(p)$ is
all-ones, $\Prob[a(p)\mid F]=\tfrac12$ and $\Prob[a(p)\mid M]=1$, while
$\Prob[a(p)\mid U]=0$: hence
\[
\Prob[p\in F\mid T_\ell]
=\frac{(f/2)\,\Theta_F}{(f/2)\,\Theta_F+m\,\Theta_M}
\;\le\;\frac{f/2}{f/2+m}\bigl(1+O(e/n)\bigr)
=q\bigl(1+O(e/n)\bigr),
\]
using $\Theta_F/\Theta_M\in(1-O(e/n),\,1+O(e/n))$.  If $a(p)$ contains
a $0$, then $\Prob[a(p)\mid M]=0$ and $\Prob[a(p)\mid U]=1$, so the
posterior is at most
$\frac{f/2}{f/2+(1-f-m)}\bigl(1+O(e/n)\bigr)\le f\bigl(1+O(e/n)\bigr)$.
In both sub-cases
$\Prob[p\in F\mid T_\ell]\le\max(q,f)(1+O(e/n))$, which is
Lemma~\ref{lem:bayes} relaxed by the perturbation factor.

Case 2: $p$ is not queried on the path.  Then $T_\ell$ reveals no
own-answer for $p$, and Lemma~\ref{lem:depletion} with $S$ the queried
cells gives $\Prob[p\in F\mid T_\ell]\le f(1+O(e/n))$ directly.

Summing over leaves and using
$\sum_\ell\Prob[\mathrm{reach}(\ell)]\le 1$ completes the proof.  For
the corollary, $q\le\frac12+o(1)$ is equivalent to $4d^2\le(1+o(1))n$,
i.e.\ $2d\le(1+o(1))\sqrt{2n}$.
\end{proof}

\begin{remark}[scope: Theorem~\ref{thm:B} is a budgeted cap]
\label{rem:Bscope}
The coupling slack $O(e/n)$ is budget-dependent, and the bound is
vacuous at $e\asymp n^2$.  The full-scan counting tree of
Theorem~\ref{thm:floor} queries $n(n+1)$ single variables and succeeds
with probability $1-2^{-\Theta(d^2)}$; this does \emph{not} contradict
Theorem~\ref{thm:B}: the counting tree lives where the slack has
exploded.  At the program's budget $e=d\log k\ll n$ the slack is
$o(1)$ and Theorem~\ref{thm:B} is the strongest proved barrier for the
adaptive single-variable class.  Reconciling the unbounded and
budgeted quantifiers is Section~\ref{sec:budget}.
\end{remark}

\begin{remark}[measured confirmations]\label{rem:measured-adaptive}
Exact channel simulations (sampling the proved channel law) confirm
Theorem~\ref{thm:B} and Lemma~\ref{lem:bayes}.
\begin{enumerate}
\item At $(n,d)=(32,2)$: the analytic posterior is
$q=\frac{2d^2+d}{2d^2-d+n}=\frac{10}{38}=0.2632$; the measured freeness
posterior given answer $1$ was $0.2367$ over $300$ simulations, against
the base rate $f=0.0189$, and the first-hit scan over $33$ queries
succeeded $0.2367$.
\item Threshold scan at $n=128$: first-hit success tracks the posterior
across $d=2$, $4$, $8$, $11$, $16$, $24$, $32$: measured success
$0.018$, $0.072$, $0.246$, $0.400$, $0.656$, $0.898$, $0.970$, against
the posterior $q=0.075$, $0.231$, $0.548$, $0.705$, $0.846$, $0.936$,
$0.970$; the posterior crosses
$\frac12$ at $d\approx\sqrt{n/2}=8$ as predicted, and the scan's own
success crosses shortly above it, the offset being the budget factor
$1-(1-p_1)^{s}$ with per-query hit probability $p_1=f/2+m$.
\item Scaling at fixed $(d,s)=(2,25)$: success
$0.1450$, $0.0425$, $0.0125$ at $n=32$, $64$, $128$ ($400$ simulations
each), an
empirical exponent of $-1.77$ against the cap model's $-2$: the
freeness signal decays polynomially in $n$ exactly as the rare-event
analysis predicts, while $k^{-O(1)}=2^{-n^{\delta}O(1)}$ decays
exponentially, so the hypothesis' margin grows rapidly in the
conjecture's regime.
\item Adaptive degree-$2$ strategies at $(d,s)=(2,40)$: single-variable
$0.183$, $0.067$, $0.013$; hybrid (single plus AND) $0.250$, $0.083$,
$0.020$; pure AND-scan $0.027$, $0.003$, below $10^{-3}$; at $n=32$,
$64$, $128$: empirical exponents
$-1.8$ to $-1.9$; no tested adaptive degree-$2$ strategy decays more
slowly than $n^{-2}$ (see Section~\ref{sec:and} and Open
Problem~\ref{op:chaining}).  (The tested ``hybrid'' branch was later
found to have degenerated to the plain scan --- its product test was
never implemented --- so this probe never exercised
adjacency-conditioned certification; the dedicated chaining study of
Section~\ref{sec:adjacency} supplies the missing test and reverses the
empirical conclusion at the tree level.)
\end{enumerate}
\end{remark}

\section{Degree-2 product queries: the no-lift theorem}\label{sec:and}

At $p=2$ we formalize the answer to a product query
$g=x_{ab}\cdot x_{cd}$ by the multiplicative evaluation semantics, as in
the session's simulation harness and consistently with the channel law:
\[
a(g)\;=\;a(x_{ab})\cdot a(x_{cd}),
\]
so $a(g)=1$ iff both factors evaluate to $1$.  (See
Remark~\ref{rem:linear-reading} for the strictly linear reading
$L(g^\rho)$; the two agree except when both component pairs are free.)

\begin{theorem}[AND-query no-lift]\label{thm:and}
In the pipeline at $p=2$ (equivalently, in the coin channel), let
$g=x_{ab}\cdot x_{cd}$ with distinct component pairs and let
$pa=(a,b)$ be the first component.  Then
\[
\Prob[a(g)=1]\;=\;\frac{f^2}{4}+fm+m^2,
\qquad
\Prob\bigl[a(g)=1 \text{ and } pa\in F\bigr]\;=\;\frac{f^2}{4}+\frac{fm}{2},
\]
and consequently the per-hit labeling posterior is
\[
\Prob\bigl[pa\in F\mid a(g)=1\bigr]
=\frac{f^2/4+fm/2}{f^2/4+fm+m^2}
=\frac{f}{f+2m}
=\frac{f/2}{f/2+m}
\;=\;q ,
\]
exactly the single-variable posterior: the AND query is a rarer joint
test whose per-hit posterior coincides with the single-variable Bayes
posterior.
\end{theorem}

\begin{proof}
Enumerate the joint statuses of the two component pairs, using that the
answers of distinct free pairs are independent fair bits
(Theorem~\ref{thm:channel}(iii) with Lemma~\ref{lem:completion}(a)) and
that matched and killed--unmatched answers are determined:
\begin{center}
\begin{tabular}{lccc}
\hline
statuses of $(pa,pb)$ & probability & $\Prob[a(g)=1]$
 & $\Prob[a(g)=1,\ pa\in F]$\\
\hline
both free & $f^2$ & $1/4$ & $1/4$\\
$pa$ free, $pb$ matched & $fm$ & $1/2$ & $1/2$\\
$pa$ matched, $pb$ free & $fm$ & $1/2$ & $0$\\
both matched & $m^2$ & $1$ & $0$\\
some killed--unmatched & $1-(f+m)^2$ & $0$ & $0$\\
\hline
\end{tabular}
\end{center}
For ``both free'' the answer is the product of two independent fair
design bits, equal to $1$ with probability $\frac14$, and $pa$ is then
free with certainty.  For ``free--matched'' the matched factor
contributes its determined $1$ and the free factor its fair bit.  Column
sums give the two displayed identities.  For the posterior, multiply
numerator and denominator by $4/f$ (a polynomial identity in $(f,m)$,
hence exact):
\[
\frac{f^2/4+fm/2}{f^2/4+fm+m^2}
=\frac{f+2m}{f+4m+4m^2/f}
=\frac{f(f+2m)}{f^2+4fm+4m^2}
=\frac{f}{f+2m}
=\frac{f/2}{f/2+m},
\]
using $(f+2m)^2=f^2+4fm+4m^2$.
\end{proof}

The matched--matched mass dominates $\Prob[a(g)=1]$ and contributes zero
labeling success, exactly canceling the both-free advantage: within the
class of trees whose degree-$2$ tests are single-pass, no lift over the
single-variable posterior is possible.  At the tree level the picture
changed with the adjacency mechanism (Section~\ref{sec:adjacency}): the
lift exists, through certainty events rather than better posteriors,
and it needs single variables, not products.  For degree-$2$ chains
specifically, budget-honest accounting shows no lift: the AND-scan is
\emph{worse} than the single-variable scan ($0.1333$ against $0.2550$
at $(32,2)$, budget $1000$), and AND-split chains that certify through
adjacency fire far more rarely than the degree-$1$ chain
(Open Problem~\ref{op:chaining}).

\begin{remark}[the strictly linear reading]\label{rem:linear-reading}
If instead the answer to $g=x_{ab}x_{cd}$ is the linear value
$L(g^\rho)$, then when both component pairs are free the answer is the
design value of the degree-$2$ monomial, itself a fair bit provided that
monomial is not a consequence of the restricted system (generic for
distinct rows and columns).  The table then changes only in the
``both free'' row, from $\frac14$ to $\frac12$, and the posterior
becomes $\frac{f(f+m)}{f^2+2fm+2m^2}$, equal to $0.311$ at $(32,2)$
against $0.263$; simultaneously the hit probability falls by the loss of
the matched--matched mass from the both-free row's complement, i.e.\ the
test becomes rarer per unit of posterior gain, and the no-lift
conclusion for first-hit strategies at fixed budget is unaffected in
order of magnitude.  Open Problem~\ref{op:chaining} subsumes both
readings.
\end{remark}

\begin{remark}[pipeline degree-$2$ laws; transfer status of
Theorem~\ref{thm:and}]\label{rem:andpipeline}
On the exact pipeline the product semantics of this section is not the
design law: the same-hole and same-pigeon degree-$2$ monomials are
consequences (determined $0$, where the product semantics answers $1$
with probability $\tfrac14$), the diagonal monomials are fresh fair
bits independent of their two component singles (the product semantics
correlates them perfectly), and the square law $L(x^2)=L(x)$ is the
one identity the product semantics gets right.  The per-hit no-lift of
Theorem~\ref{thm:and} is derived under the product semantics; its
pipeline transfer is not covered by the proof but is
measured-consistent: true monomial queries on the exact pipeline land
on the cap, $0.2629$ $[0.2129,0.3199]$ against $q=0.2632$ at $(32,2)$,
budget $1000$ (the \texttt{mono2} row of the chaining study).  The
transfer is listed as coin-model-only in
Table~\ref{tab:faithfulness}.
\end{remark}

\section{The two-phase certification tree}\label{sec:twophase}

\subsection{The tree and its exact error law in the coin channel}

The two-phase certification tree $\mathcal{T}_{2\mathrm{ph}}$ is the
adaptive degree-$1$ tree:
\begin{enumerate}
\item[Phase 1.] For pigeons $i=1,\dots,n+1$: query the row sum
$Q_i=1+\sum_jx_{ij}$; stop at the first answer $1$ and set $i^{\ast}$ to
that pigeon.  (In the coin channel an assigned pigeon answers the
determined $0$; a free pigeon answers $1$ plus the XOR of its $2d$
free-row design bits, a fair coin: answer $1$ iff the row's free count
is \emph{even}.)
\item[Phase 2.] For holes $j$: query $x_{i^{\ast}j}$; stop at the first
answer $1$ and label the leaf $(i^{\ast},j)$.
\item[Fallback.] If all queries of a phase answer $0$, label an
arbitrary pair.
\end{enumerate}

\begin{theorem}[literal error law of the two-phase tree, coin channel]
\label{thm:twophase}
In $\mathsf{Coin}(n,d)$, under the literal query semantics (the answer
to $Q_i$ is $1$ iff the row XOR is $0$), the two-phase tree errs with
probability exactly
\[
\mathrm{err}(\mathcal{T}_{2\mathrm{ph}})
\;=\;2^{-(2d+1)}\;+\;\bigl(1-2^{-(2d+1)}\bigr)\;2^{-2d+1}
\;=\;2^{-2d}\bigl(\tfrac52-2^{-2d}\bigr)
\;=\;5\cdot 2^{-(2d+1)}\bigl(1+o(1)\bigr),
\]
up to the fallback's base-rate luck: with the fallback pair free with
probability $f$, the total error is
$(1-f)\,\mathrm{err}(\mathcal{T}_{2\mathrm{ph}})$.  The first term is
the phase-$1$ miss (every free pigeon's row answers even); the second
is the \emph{count-$0$ branch}: the certified pigeon's row has an even
count of $1$s, and all-zero is even, an event of conditional
probability $2^{-2d}/\tfrac12=2^{-2d+1}$.
\end{theorem}

\begin{proof}
Phase 1 queries the row sums in order.  An assigned pigeon answers the
determined $0$: its row XOR is $1$ (the matched cell), so
$Q_i^{\rho}=1+1=0$.  A free pigeon's answer is $1$ plus the XOR of its
$2d$ free-row design bits: an independent fair coin, the events across
the $2d+1$ free pigeons being independent (disjoint rows, independent
coordinates).  An answer $1$ certifies freeness soundly (assigned
pigeons never produce it), and phase 1 misses with probability exactly
$2^{-(2d+1)}$.

Given certification, the certified pigeon's row has an even count.
Phase 2 scans that row (killed holes answer the determined $0$) and
succeeds iff the count is at least $1$; conditionally on evenness the
row's answer pattern is uniform over the $2^{2d-1}$ even patterns, of
which exactly one is all-zero, so the failure probability is
$2^{-2d}/\frac12=2^{-2d+1}$.  The earlier
draft asserted this step deterministic (``phase 2 never fails given
certification''); that assertion is precisely the gap corrected here:
an even row count need not contain a $1$.

Total, with the fallback firing on either failure and its label free
with the base rate $f$:
\[
\mathrm{err}(\mathcal{T}_{2\mathrm{ph}})
=2^{-(2d+1)}+\bigl(1-2^{-(2d+1)}\bigr)2^{1-2d}
=2^{-2d}\bigl(\tfrac52-2^{-2d}\bigr)
=5\cdot 2^{-(2d+1)}\bigl(1+o(1)\bigr). \qedhere
\]
\end{proof}

\begin{remark}[the simulator artifact; retraction of the exact law
$1-(1-f)2^{-(2d+1)}$]\label{rem:artifact}
The session's simulators certified a pigeon as free when its row XOR
equals $1$ (odd count of $1$s), after pre-checking, with oracle
knowledge of $\rho$, that the pigeon is free.  No legitimate
degree-$1$ query implements that certification: by linearity (the
mechanism of Lemma~\ref{lem:degonespan}; explicitly
Lemma~\ref{lem:linear} below) the stated query $Q_i=1+\sum_jx_{ij}$
answers $1$ iff the row XOR is $0$, the constant-free row sum answers
$1$ on \emph{every} killed pigeon, and the oracle pre-check conditions
on the very quantity (freedom) that the tree is supposed to be
certifying.  The exact law $1-(1-f)2^{-(2d+1)}$, with its ``phase $2$
never fails'' step, is accordingly \textbf{retracted}.  The previously
recorded measurements --- success $0.970$ at $(32,2)$, $0.97$--$0.98$
at $(64,2)$, $0.974$ at $(128,2)$, $0.998$ at $(64,4)$, and the
five-point replication over $816{,}000$ simulations --- verify the
\emph{hybrid} law of the artifact simulator (certify on ``free and
row-XOR odd''), not the tree; the replication's $(48,4)$ fluctuation
investigation stands, but within the artifact.
Table~\ref{tab:twophase} separates the hybrid rows from the literal
rows, the latter confirming Theorem~\ref{thm:twophase} point for
point, including by exact enumeration.
\end{remark}

\begin{table}[t]
\centering
\small
\begin{tabular}{llrrr}
\hline
$(n,d)$ & law & measured & predicted & simulations\\
\hline
\multicolumn{5}{l}{\emph{hybrid-artifact rows: certify on ``free and
row-XOR odd'' (Remark~\ref{rem:artifact})}}\\
$(32,2)$ & $1-(1-f)2^{-5}$ & 0.970 & 0.969 & 500\\
$(64,2)$ & $1-(1-f)2^{-5}$ & 0.97--0.98 & 0.969 & 500\\
$(128,2)$ & $1-(1-f)2^{-5}$ & 0.974 & 0.969 & 500\\
$(64,4)$ & $1-(1-f)2^{-9}$ & 0.998 & 0.998 & 500\\
\multicolumn{5}{l}{\emph{literal-tree rows: the law of
Theorem~\ref{thm:twophase}}}\\
$(32,2)$ & $1-(1-f)\bigl[2^{-5}+(1-2^{-5})2^{-3}\bigr]$
& 0.8500 & 0.8505 & 2000\\
$(64,2)$ & $1-(1-f)\bigl[2^{-5}+(1-2^{-5})2^{-3}\bigr]$
& 0.8380 & 0.8484 & 1500\\
$(64,4)$ & $1-(1-f)\bigl[2^{-9}+(1-2^{-9})2^{-7}\bigr]$
& 0.9888 & 0.9904 & 800\\
$(32,2)$ & retry: $1-(1-f)(9/16)^{5}$ & 0.9430 & 0.9448 & 2000\\
\hline
\end{tabular}
\caption{Measured success of the two-phase certification tree,
separated by law.  Hybrid rows: the artifact simulator of
Remark~\ref{rem:artifact}, whose certification condition no degree-$1$
query implements; they verify the hybrid law, not the tree.  Literal
rows: an exact literal-law re-implementation confirming
Theorem~\ref{thm:twophase}.  Exact enumeration at $(5,2)$ (all
$2^{20}$ coin matrices) gives strict values $0.8476562=(31/32)(7/8)$
for the literal tree and $0.9436865=1-(9/16)^5$ for the literal retry,
both digit-exact; the published-law tree re-measures $0.9695$ at
$(32,2)$, reproducing the corpus's $0.97$ (flipped-law audit only).}
\label{tab:twophase}
\end{table}

\begin{remark}[retry variant, literal law]\label{rem:retry}
If phase 1 collects the full certified set $C$ (the free pigeons whose
row answered even, $|C|\sim\mathrm{Bin}(2d+1,\tfrac12)$) and phase 2
retries across $C$, the rows being disjoint, the failure event is
``every certified row all-zero'', of probability
\[
\Bigl(\tfrac{1+2^{1-2d}}{2}\Bigr)^{2d+1}
=2^{-(2d+1)}\bigl(1+o(1)\bigr):
\]
the retry recovers the retracted law's asymptotic constant.  Measured
(literal re-implementation): success $0.9430$ at $(32,2)$ against
$0.9448$ predicted; exact enumeration at $(5,2)$: strict success
$0.9436865=1-(9/16)^5$, digit-exact.  The earlier retry measurements
($0.975/0.968/0.970$ at $(n,d)=(32,2),(64,2),(128,2)$, $0.9925$ and
$1.0000$ at $d=4$, $1.0000$ at $d=8$) were produced by the artifact
simulator and verify the hybrid law (Remark~\ref{rem:artifact}).
\end{remark}

\subsection{Model fidelity: the rigidity against the literal pipeline}

Theorem~\ref{thm:twophase} is a theorem about the coin channel.  Against
the literal pipeline of Definition~\ref{def:pipeline} the phase-1
certification channel does not exist.

\begin{corollary}[the parity-certification channel is closed against
the pipeline]\label{cor:closed}
Against the literal pipeline of Definition~\ref{def:pipeline} the
two-phase tree's phase 1 never certifies: every pigeon answers the
determined $0$ to its row sum (Proposition~\ref{prop:rigidity}), phase
1 completes without firing, and the fallback decides, so the tree's
success collapses to the base rate $f$, far below either coin-channel
value of Theorem~\ref{thm:twophase} (hybrid or literal).  Measured on
the exact pipeline (uniform designs over $\Des(4,2)$) at $(32,2)$,
budget $1000$: two-phase $0.2767$ $[0.2323,0.3260]$ against the
single-variable scan's $0.2650$ $[0.2213,0.3138]$ ---
statistically indistinguishable, the phase-$1$ signal void ($0$
answers $1$ among $2000$ free-pigeon row-sum queries).  The
``reversed tree'' scenario of the previous draft (stop at the first
$0$, then scan the certified row for a $1$, winning with probability
exactly $1$) does not exist: it required assigned pigeons to answer
$1$, which the zeroing restriction semantics excludes under either
presentation (Remark~\ref{rem:onekernel}); every pigeon answers $0$,
so a first $0$ certifies nothing.  The parity-certificate mechanism is
dead at kernel level; the certificate mechanisms that survive against
the pipeline are the counting certificates
(Lemma~\ref{lem:certificates}, over budget) and the adjacency
certificates (Section~\ref{sec:adjacency}).
\end{corollary}

\begin{proof}
By Proposition~\ref{prop:rigidity}, $Q_i^{\rho}$ is a restricted
consequence for free pigeons and the zero polynomial for assigned
pigeons, so every row-sum answer is the determined $0$ and phase 1
never fires; the fallback labels an arbitrary pair, free with
probability $f$.  For the reversed scenario: the first $0$ is answered
by every pigeon, free or assigned, so it certifies nothing; with the
assigned-row variables zeroed there is no reading under which an
assigned pigeon answers $1$.  The measured values are the exact-pipeline
rows of the chaining study (Section~\ref{sec:adjacency}).
\end{proof}

\begin{remark}[what the fidelity means]\label{rem:fidelity}
Three lessons.  First, the hybrid rows of Table~\ref{tab:twophase}
validate the artifact simulator's certification law, and the literal
rows validate Theorem~\ref{thm:twophase}: a
sampler that draws free-pair design bits independently implements the
coin channel, not the design space, whose row parities are forced
(Proposition~\ref{prop:rigidity}); the coin channel drops exactly the
row-parity constraints (legal fraction $\approx2^{-(2d+1)}$,
Remark~\ref{rem:onekernel}) and, at degree $2$, the product identities
(Remark~\ref{rem:andpipeline}).  Second, the rigidity is good for
the pseudo-solution program: the literal $\Omega(n,d)$ at $p=2$ is
strictly more robust against parity-certification attacks than the
coin model suggests, and both the solution condition
(Definition~\ref{def:solution}) and the chi-task hypothesis
(Theorem~\ref{thm:chi}) survive the strongest known adaptive attacker
(Section~\ref{sec:adjacency}).  Third, aggregate queries are not fully
dead in principle: a full scan of a row reveals its parity, and an
assigned row always shows exactly one answer $1$ while a free row
shows an odd number of them (parity-locking,
Theorem~\ref{thm:channel}(iv)); at $d=2$ a free row shows three $1$s
with probability $\frac12$ (odd counts $1$ and $3$ being equally
likely, $\binom{4}{c}/2^{3}$ on odd $c$), a detectable event.  But each
full scan costs $2d$ queries against a free-row density
$\approx 2d/n$, so within budget $e=d\log k$ the detection probability
is $O(ed/n^2)$-small, consistent with the cap of Theorem~\ref{thm:B}.
Quantifying the best statistical freeness signal obtainable from
aggregate queries within budget is Open Problem~\ref{op:floor}.
\end{remark}

\section{The certification floor at degree 1}\label{sec:floor}

This section settles the unbounded-budget certification problem at
degree $1$.  All trees below query single outer variables only and have
no query budget; adaptivity is allowed throughout but will not help.

\begin{definition}[answer table; configuration]\label{def:table}
The \emph{answer table} is $M\in\F_2^{(n+1)\times n}$ with
$M(i,j)=\mathrm{ans}(x_{ij})$; call $M(i,j)=1$ a \emph{one-entry}.  A
\emph{configuration} is a pair $(Z,c)=\bigl((D,R,\mu),\text{coins}\bigr)$
of the pipeline: the status data $Z$ (the free sets $D,R$ and the
matching $\mu$) together with the free-region coin draws $c$; $M$ is a
function of the configuration.  In this section the free cells'
answers are the independent fair bits of $\mathsf{Coin}(n,d)$
(Definition~\ref{def:coin}), the answer-table model that the session's
exact-channel samplers implement and in which the closed forms below
are verified; fidelity to the literal pipeline is discussed in
Remark~\ref{rem:floorfidelity}.  Fallback convention: a tree that
exhausts its search outputs a fixed fallback pair, free with the base
rate $f=(2d+1)2d/((n+1)n)$; ``strict'' exact values below carry no
fallback luck (the canonical fallback pair is matched).
\end{definition}

\begin{lemma}[linearity of answers]\label{lem:linear}
For any affine-linear query $g=a+\sum_{p\in S}x_p$ ($a\in\F_2$, $S$ any
set of pairs: the entire degree-$1$ query class),
$\mathrm{ans}(g)=a+\bigoplus_{p\in S}\mathrm{ans}(x_p)$.
\end{lemma}

\begin{proof}
$\mathrm{ans}(g)=L(g^{\rho})
=L\bigl(a\cdot 1+\sum_px_p^{\rho}\bigr)
=a\,L(1)+\sum_pL(x_p^{\rho})
=a+\sum_p\mathrm{ans}(x_p)$ over $\F_2$, using the linearity of $L$
and the design axiom $L(1)=1$.
\end{proof}

\begin{corollary}[full-scan dominance]\label{cor:fullscan}
For every adaptive degree-$1$ tree $T$ there is a \emph{non-adaptive}
tree $T'$ that queries all $n(n+1)$ single variables and then outputs,
with the same success probability as $T$ on every configuration.
Consequently the optimal error over the whole adaptive degree-$1$
family equals the Bayes error of the pair posterior given $M$,
\[
\mathrm{err}^*\;=\;\mathbb{E}_M\Bigl[\,1-\max_p\Prob\bigl[p\in F\mid
M\bigr]\Bigr],
\]
and randomized trees do not beat the deterministic optimum (convexity).
\end{corollary}

\begin{proof}
By Lemma~\ref{lem:linear} the answer to any query of $T$ is an
$\F_2$-function of $M$.  By induction on the query index, $T$'s
$t$-th query and its answer are functions of $M$ for every $t$, so
$T$'s output label is a function of $M$.  $T'$ computes $M$ by a fixed
full scan, then replays $T$'s decision procedure, answering each
simulated query by Lemma~\ref{lem:linear}, and outputs $T$'s final
label.  The displayed identity is the definition of the optimal
measurable selector of a pair given $M$; a randomized tree is a
mixture of deterministic trees, so the optimum is attained by a
deterministic one.
\end{proof}

\begin{lemma}[counting certificates]\label{lem:certificates}
In every configuration, every killed pigeon's row of $M$ contains
exactly one one-entry (at its matched hole) and every killed hole's
column contains exactly one one-entry (at its matched pigeon).  Hence,
soundly against every configuration consistent with the transcript:
\begin{enumerate}
\item[(a)] a row containing any number of $1$s other than exactly one
certifies its pigeon free;
\item[(b)] a column containing any number of $1$s other than exactly
one certifies its hole free;
\item[(c)] a one-entry in the row of a certified-free pigeon certifies
its hole free, and a one-entry in the column of a certified-free hole
certifies its pigeon free.
\end{enumerate}
\end{lemma}

\begin{proof}
Killed pigeon $i$: $x_{i,\mu(i)}$ answers $1$ (matched,
Theorem~\ref{thm:channel}(i)); every other pair $(i,j')$ of the row is
killed--unmatched and answers $0$ (Theorem~\ref{thm:channel}(ii));
there is no free hole in the row because $i\notin D$.  Symmetrically
for columns.  Soundness of (a): if the pigeon were killed, its row
would show exactly one $1$.  Soundness of (c): with the pigeon
certified free, a one-entry $(i,j)$ cannot be matched, and killed $j$
would make $(i,j)$ killed--unmatched (answer $0$).
\end{proof}

The session's earlier trees aggregate answers by XOR
($Q_i=1+\text{row XOR}$, Theorem~\ref{thm:twophase}).  The certificates
above need the \emph{count}, which a decision tree may compute but no
single parity query can express; this is the whole gap between
Theorem~\ref{thm:twophase} and the floor, and the mechanism that the
parity framing missed.

\begin{lemma}[pigeonhole rescue]\label{lem:rescue}
If every row of $M$ shows exactly one $1$, then some column shows at
least two $1$s (there are $n+1$ rows and $n$ columns).  In that event
the column is certified free by Lemma~\ref{lem:certificates}(b), a
certified-free column contains no matched pair, so every one-entry in
it certifies its pigeon free by Lemma~\ref{lem:certificates}(c), and
the tree outputs a \emph{certified} free pair (error $0$ on this
event).
\end{lemma}

\begin{lemma}[the near-permutation class]\label{lem:classb}
Define class $\mathfrak{b}=\{$every column of $M$ shows exactly one
$1$, and every row shows at most one $1\}$.  Then:
\begin{enumerate}
\item[(i)] $\Prob[\mathfrak{b}]=(2d+1)\,(2d)!\,/\,2^{(2d+1)2d}$;
\item[(ii)] on class $\mathfrak{b}$ the free-region coin matrix is a
permutation matrix on $2d$ of the $2d+1$ free pigeons plus one
all-zero row, and every one-entry of $M$ lies in the near-permutation
formed by the killed matching and that partial permutation;
\item[(iii)] the configurations consistent with a class-$\mathfrak{b}$
table are exactly the $\binom{n}{n-2d}$ choices of $n-2d$ one-entries
as the matched pairs; each hole is free in a fraction exactly $2d/n$
of them, the all-zero row's pigeon is free in all of them, and no pair
has posterior above $2d/n$;
\item[(iv)] the Bayes error on class $\mathfrak{b}$ is exactly
$1-2d/n$.
\end{enumerate}
\end{lemma}

\begin{proof}
(i) Killed rows and columns show exactly one $1$ automatically, so the
condition lives on the $(2d+1)\times 2d$ free-region coin matrix: each
of the $2d$ columns has exactly one $1$, and the rows sum to $2d$ with
each row at most one $1$, forcing exactly one all-zero row and a
bijection elsewhere: $(2d+1)$ choices times $(2d)!$ permutation
matrices, out of $2^{(2d+1)2d}$ equally likely coin matrices.
(ii) Immediate.  (iii) A consistent configuration is any choice of
$n-2d$ one-entries as $\mu$, with killed sets the endpoints of the
chosen entries: disjoint rows and columns hold automatically (each
column has one one-entry, rows at most one), every unchosen one-entry
has both endpoints free, and the coins are forced; conversely every
configuration generating this $M$ arises so.  Each hole's column edge
is chosen in $\binom{n-1}{n-2d-1}$ of the $\binom{n}{n-2d}$ twins, so
$\Prob[\text{hole free}]
=1-\binom{n-1}{n-2d-1}\big/\binom{n}{n-2d}=2d/n$; the all-zero row
cannot be matched, so its pigeon has posterior $1$; every
matched-candidate pigeon has posterior $1-2d/n$; and hence every pair
has joint posterior at most $2d/n$ (pairs with the all-zero row:
exactly $2d/n$; all other pairs: bounded by the hole's $2d/n$).
(iv) Immediate.
\end{proof}

\begin{theorem}[exact certification floor at degree $1$]\label{thm:floor}
Every adaptive degree-$1$ tree for the pipeline at $p=2$, with no
query budget, in the answer-table model of Definition~\ref{def:table},
has error probability exactly
\[
\mathrm{err}^*\;=\;\Bigl(1-\frac{2d}{n}\Bigr)\cdot
\frac{(2d+1)\,(2d)!}{2^{(2d+1)2d}}
\;=\;2^{-4d^2+2d\log_2 d-O(d)}
\;=\;2^{-\Theta(d^2)},
\]
and the following \emph{non-adaptive} tree attains it: query all
$n(n+1)$ variables in fixed order; compute all row and column counts;
then
\begin{enumerate}
\item[(1)] if some row and some column both have count $\neq1$, output
(that row's pigeon, that column's hole);
\item[(2)] else if every row has count exactly $1$, output the first
one-entry of the first column with count $\ge2$;
\item[(3)] else if every column has count exactly $1$ and some row has
count $\ge2$, output that pigeon with its first $1$-column;
\item[(4)] else (class $\mathfrak{b}$) output the count-$0$ row's
pigeon with any hole.
\end{enumerate}
\end{theorem}

\begin{proof}
The four cases exhaust all $M$.  Cases 1--3 output pairs with
posterior $1$ by Lemmas~\ref{lem:certificates} and~\ref{lem:rescue},
the best possible.  Case 4 is class $\mathfrak{b}$, where by
Lemma~\ref{lem:classb} the posterior of every pair is at most $2d/n$
and the tree attains it.  By Corollary~\ref{cor:fullscan} no tree
exceeds the Bayes value on any $M$, so the displayed value is the
exact optimum.  The asymptotic form is Stirling's applied to
$\log_2\mathrm{err}^*
=-(2d+1)\,2d+\log_2\bigl((2d+1)(2d)!\bigr)+O(1)$.
\end{proof}

At $(32,2)$: $\mathrm{err}^*=1.0\times10^{-4}$, against the conjectured
floor $2^{-5}=3.1\times10^{-2}$ of the earlier draft and
Theorem~\ref{thm:twophase}'s literal error $39/256=1.5\times10^{-1}$;
at $(64,4)$: $\mathrm{err}^*=6.7\times10^{-17}$ against $2^{-9}$.
The $2^{-\Theta(d)}$ certification-floor conjecture and the two-phase
tree's optimality are both false.

\begin{remark}[verification record]\label{rem:floorverify}
All closed forms are verified digit-exactly by exhaustive computation
in one registered run of \texttt{cert\_floor\_check.py} ($31$\,s; fixed
seeds, registered output reproduced in the corpus).
\begin{enumerate}
\item \emph{Bayes audit at $(4,1)$, exhaustive.}  All
$120\times64=7680$ configurations enumerated and grouped into $6120$
answer-table classes.  On every class the counting tree's output pair
achieves exactly the maximum pair posterior (max
$\lvert\text{tree}-\text{best}\rvert=0$), and the pooled Bayes error
$\mathbb{E}_M[1-\max\text{ posterior}]=0.046875$ equals
$\mathrm{err}^*(4,1)$ exactly.
\item \emph{Exact enumeration at $(5,2)$.}  All $2^{20}$ coin
matrices: strict success $0.9998856=1-120/2^{20}$, matching
$\Prob[\mathfrak{b}]=(2d+1)(2d)!/2^{(2d+1)2d}$ digit for digit.
\item \emph{Monte Carlo.}  Counting-tree success $1.0000$ at $(32,2)$
in $2000$ runs (expected failures $0.2$ at
$\mathrm{err}^*=1.0\times10^{-4}$), $1.0000$ at $(64,2)$ in $1500$
runs, and $1.0000$ at $(64,4)$ in $800$ runs
($\mathrm{err}^*=6.7\times10^{-17}$).
\end{enumerate}
In a $200{,}000$-configuration sample at $(5,2)$, seven
multi-configuration classes showed a nonzero tree-versus-Bayes gap;
every one is a class-$\mathfrak{b}$ table, where
Lemma~\ref{lem:classb}(iii) makes all pair posteriors flat at $2d/n$
and small samples fluctuate.  A gap on any other table would have
refuted Theorem~\ref{thm:floor}; none occurred.
\end{remark}

\begin{proposition}[rows-and-columns parity tree, $O(n)$ queries]
\label{prop:rcparity}
Scan the $n+1$ row sums $Q_i$ and the $n$ column duals
$K_j=1+\sum_ix_{ij}$ ($2n+1$ queries); let $I$ be the set of pigeons
that answered $1$ and $J$ the set of holes that answered $1$ (both
certified free, soundly: assigned lines answer the determined $0$).
Output $(\min I,\min J)$ if both are nonempty; if $I$ is empty, the
global parity identity (the sum of the row XORs equals the sum of the
column XORs) forces some column to be certified, and the tree scans
the certified columns for a one-entry and outputs it; symmetrically if
$J$ is empty.  Its error is exactly
\[
\mathrm{err}_{R}
=\frac{\displaystyle\sum_{b\ \mathrm{odd}}\binom{2d}{b}\,
2^{2d(2d-b-1)}
+\sum_{a\ \mathrm{odd}}\binom{2d+1}{a}\,
2^{(2d-a)(2d-1)}}{2^{(2d+1)2d}}
=(6d+2)\,2^{-6d}\,\bigl(1+o(1)\bigr),
\]
versus Theorem~\ref{thm:twophase}'s $5\cdot 2^{-(2d+1)}$: an
exponential gap of $2^{4d}/\mathrm{poly}(d)$ at equal query cost.  The
purely non-adaptive variant (output $(\min I,\min J)$ only when both
sets are nonempty, fallback otherwise) has error $2^{-(2d+1)}+2^{-2d}$.
Exact enumeration at $(5,2)$: strict success
$0.9965019=1-3668/2^{20}$ (the two failure forms contributing $1028$
and $2640$), digit-exact; measured $0.9955$ at $(32,2)$, $0.9960$ at
$(64,2)$, $1.0000$ at $(64,4)$.
\end{proposition}

\begin{proof}
Certification is sound: assigned lines answer the determined $0$, so
an answer $1$ certifies freeness.  The non-adaptive variant errs iff
no row or no column is certified, events of probability $2^{-(2d+1)}$
and $2^{-2d}$; they are disjoint, because both failing forces the sum
of the row XORs and the sum of the column XORs to differ ($n+1$ odd
versus $n$ odd), contradicting the global parity identity.  For the
adaptive tree, the rescue branches leave two failure forms, F1 (no row
certified and the column rescue failing:
$\sum_{b\ \mathrm{odd}}\binom{2d}{b}2^{2d(2d-b-1)}$ configurations) and
F2 (rows certified but no column and the row rescue failing:
$\sum_{a\ \mathrm{odd}}\binom{2d+1}{a}2^{(2d-a)(2d-1)}$); dividing by
$2^{(2d+1)2d}$ and applying Stirling gives the displayed asymptotics.
The case analysis is mechanical; the closed form is confirmed
digit-exactly by the exact enumeration quoted.
\end{proof}

\begin{remark}[floor fidelity: coin model versus literal pipeline]
\label{rem:floorfidelity}
Theorem~\ref{thm:floor} is proved and verified in the answer-table
model of Definition~\ref{def:table}, in which the free cells' answers
are independent fair bits.  The full scan on which the counting tree
relies leaves the regime of Corollary~\ref{cor:coin}: it pins every
row in its entirety, the event $E_{\mathrm{full}}$ on which the
coin-channel equivalence for single-variable trees fails, and
Proposition~\ref{prop:rigidity}'s forced row parities then constrain
the literal pipeline's answer table (no free row of the literal
pipeline is ever all-zero).  A verbatim transfer of the floor to the
literal pipeline's answer table is therefore not claimed here; the
chi-transfer note of the corpus (for the counting tree the
chi-indicator coincides with its error indicator whenever no
degree-$\le d$ consequence in $V(n,d)^{\rho}$ has constant term $1$, a
reading of Definition~\ref{def:pipeline} to verify against
\cite{Kra26}) transfers the floor to the chi-task directly under that
reading.  The budgeted floor, which is what
Theorem~\ref{thm:chi} actually quantifies over, is
Open Problem~\ref{op:floor}.
\end{remark}

\begin{remark}[Theorem~\ref{thm:floor} under parity-locking]
\label{rem:parityfloor}
The exact constant
$\mathrm{err}^*=(1-2d/n)(2d+1)(2d)!/2^{(2d+1)2d}$ is a
\emph{coin-model} number: it is proved and verified in the
answer-table model of Definition~\ref{def:table}, whose free rows are
unconstrained.  On the literal pipeline every free row is
parity-locked (Theorem~\ref{thm:channel}(iv)): a fully scanned free
row shows an odd count of $1$s, with
$\Prob[\text{count}=c]=\binom{2d}{c}/2^{2d-1}$ on odd $c$ --- so the
count-$1$ probability is $2d/2^{2d-1}$, exactly twice the coin model's
$2d/2^{2d}$ --- while killed rows still show exactly one $1$.  The
counting mechanism survives (a count other than $1$ still certifies
freeness soundly, Lemma~\ref{lem:certificates}), but the hard class
does not merely shift: it \emph{vanishes}.  Class
$\mathfrak{b}$ (Lemma~\ref{lem:classb}) would require an all-zero free
row: with $2d+1$ free rows and $2d$ free holes each showing exactly
one $1$, only $2d$ rows can carry a $1$ --- and parity-locking
(forced row XOR $1$) forbids exactly the all-zero row that
$\mathfrak{b}$ needs.  Hence on the literal pipeline the full-scan
counting tree succeeds \emph{with probability exactly $1$}: the
unbounded floor terminates at zero, not at a changed constant
(verified: $0$ failures in $5000$ exact-channel simulations).
Theorem~\ref{thm:floor} therefore survives as the exact floor of the
\emph{coin channel} --- the row-incomplete completion model, which is
the correct model of the pipeline below per-row query cost
$\Theta(n)$ --- and the open remnant is the \emph{budgeted}
all-degrees floor (Open Problem~\ref{op:floor}), whose
degree-$\le 2$ form is settled at
$\mathrm{err}^*(d,\,d\log k)=k^{-\Theta(d^2/n)}$.
\end{remark}

\section{Adjacency certificates and the confirmation chain}
\label{sec:adjacency}

The caps proved so far bound what a single answer can say: the per-hit
posterior of any scan hit is $q$ (Lemma~\ref{lem:bayes},
Theorem~\ref{thm:and}), and the measured confirmations show every
single-hit strategy landing on $q$.  This section exhibits the
mechanism that defeats the per-query caps at the tree level.  It rests
on a certificate that consumes only the determined answers
(Theorem~\ref{thm:channel}(i)--(ii)), so it is sound against the
literal pipeline of Definition~\ref{def:pipeline} directly, with no
coin-channel intermediate step and no property of the design values at
all.

\begin{lemma}[adjacency certification]\label{lem:adjacency}
Let $p\neq p'$ be pairs sharing a row or a column.  If
$\mathrm{ans}(x_p)=\mathrm{ans}(x_{p'})=1$, then $p$ and $p'$ are both
free.  The certificate is error-free, and it is valid against the
literal pipeline (Remark~\ref{rem:onekernel}'s single design space).
\end{lemma}

\begin{proof}
An answer $1$ excludes the killed--unmatched status, which answers the
determined $0$ (Theorem~\ref{thm:channel}(ii)); so each of $p,p'$ is
free or matched.  A matched pair forces all of its row-neighbors and
column-neighbors to the determined $0$: if $\rho(i)=j$, then every
other cell of row $i$ is killed--unmatched, and so is every other cell
of column $j$ (the hole $j$ is taken, and $\rho$ is injective).
Suppose $p'=(i,j)$ were matched; then $p$, being one of its
row- or column-neighbors, would answer the determined $0$ ---
contradicting $\mathrm{ans}(x_p)=1$.  So $p'$ is free.  Symmetrically,
if $p$ were matched, its neighbor $p'$ would be killed--unmatched and
answer $0$; and $p$ itself is not killed--unmatched; hence $p$ is
free.  Every step is a $\rho$-determined answer law, which no design
can alter.
\end{proof}

\subsection{The confirmation chain and its measurements}

The \emph{confirmation chain} (\texttt{confirm} in the session's
chaining study) spends a budget $B$ of single-variable reads as
follows: scan fresh pairs in random order; on a hit at $p$ (answer
$1$), spend up to $k=6$ reads on fresh row- and column-neighbors of
$p$; any neighbor answering $1$ certifies $p$ free by
Lemma~\ref{lem:adjacency} and is output immediately; otherwise resume
scanning; at budget end, output the least-confirmed accumulated hit.
The certificate events are what the caps cannot price:
Lemma~\ref{lem:adjacency} makes their posterior exactly $1$, so the
chain's success is the probability of \emph{reaching} a certainty
event, a quantity that grows with budget, whereas every previously
capped quantity was budget-independent.

\begin{remark}[measured chain performance]\label{rem:chainmeasured}
Registered run (\texttt{chi\_and\_chain.py}: $600$ simulations per
cell, Wilson $99\%$ confidence intervals, budget-honest accounting;
the self-check ``every certified output is free'' was asserted per
simulation and never fired in $52{,}800$ simulations).
\begin{enumerate}
\item \emph{Stipulated coin channel.}  Success $0.9317$
$[0.9001,0.9538]$ at budget $B=1000$ at \emph{both} $(32,2)$ and
$(64,4)$, against the single-variable scan's $0.2550$
$[0.2120,0.3033]$ and $0.3800$ $[0.3305,0.4321]$ (a $3.7\times$ lift
at $(32,2)$); the confidence intervals are disjoint from the scan's at
every $B\ge300$ at both sizes and overlap at $B=100$.  Posterior given
certification $1.0000$ in every cell.  The certificate rate is
exponential in budget: $5.17$, $20.50$, $49.17$, $91.83\%$ at
$B=30,100,300,1000$ and $(32,2)$; $7.50$, $20.83$, $50.17$,
$91.83\%$ at $(64,4)$; i.e.\ $-\ln(1-\mathrm{cert})/B$ lies in
$[0.0018,0.0026]$, roughly constant at $\approx0.0025$ per query, so
$\mathrm{cert}\approx1-e^{-B/400}$.
\item \emph{Exact pipeline} (uniform designs over $\Des(4,2)$, the
one-kernel space of Remark~\ref{rem:onekernel}): success $0.9417$
$[0.9119,0.9618]$ at $(32,2)$, $B=1000$; certificate rate $92.83\%$;
posterior given certification $1.0000$ (asserted).  The chain consumes
only the $\rho$-determined facts (matched restricts to $1$, killed to
$0$) and the marginal fairness of the free design bits (measured
$0.4918$ over $n=8000$ samples), so nothing in its mechanism changes
under parity-locking (odd row parities permit the neighbor $1$s); it
is the only tested strategy that beats the single-variable cap on the
true pipeline at $(32,2)$.
\item \emph{The per-hit no-lift is intact.}  The chain's posterior
given a plain scan hit falls \emph{below} $q$, $0.1633$
$[0.0696,0.3372]$: free hits self-select into certifications,
leaving a matched-enriched residue.  The lift is not a better guess
but a certainty reached more often; the single-pass no-lift
(Theorem~\ref{thm:and}) is sharp, and the cap it proves is a
per-query, not a per-tree, bound.
\item \emph{No lift without adjacency.}  Budget-honest AND-scans are
worse than single variables ($0.1333$ against $0.2550$ at $(32,2)$,
$B=1000$), and AND-pool and AND-split strategies sit at the cap
($0.2317$ and $0.2717$, overlapping the scan's interval); a
posterior-weighted chain that pools certified neighbors with plain
hits lifts by the same mechanism, diluted ($0.3483$ at $(32,2)$ and
$0.8000$ at $(64,4)$, $B=1000$).  The earlier ``no lift found at toy
scale'' reading is superseded: the probe that supported it never
implemented its product-test branch (a measurement bug recorded in
the corpus).
\end{enumerate}
\end{remark}

\begin{remark}[compatibility with the proved caps]\label{rem:chainvsB}
Lemma~\ref{lem:adjacency} does not contradict Theorem~\ref{thm:B} (or
the no-lift theorem): the coupling slack of the cap is
budget-dependent and vacuous at $e\asymp10^3$ ($e/n\approx31$ at
$(32,2)$, $B=1000$), while at the program's budget $e'=d\log k\ll n$
the slack is $o(1)$.  What the chain changes is the shape of the
budgeted question: its error decays in \emph{budget} like
$e^{-B/400}$ rather than staying at $q$, so the printed chi-hypothesis
survives or falls with the scaling of the certificate constant in
$(n,d)$ (Open Problem~\ref{op:adjacency}).
\end{remark}

\begin{remark}[place in the certificate trichotomy]\label{rem:trichotomy}
The known freeness certificates divide three ways.  \emph{Parity}
certificates --- query the row sum, certify on answer $1$ --- are dead
against the literal pipeline: $Q_i$ answers the determined $0$ on
every pigeon (Proposition~\ref{prop:rigidity}; measured $0/2000$ on
free pigeons), and the two-phase tree collapses to the base rate
(Corollary~\ref{cor:closed}).  \emph{Counting} certificates --- line
counts $\neq1$, Lemma~\ref{lem:certificates} --- are alive and
Bayes-optimal in the coin model unbounded (Theorem~\ref{thm:floor}),
but inert within any polylog budget (Section~\ref{sec:budget}), and
their exact constant is a coin-model number
(Remark~\ref{rem:parityfloor}).  \emph{Adjacency} certificates
(Lemma~\ref{lem:adjacency}) are alive, error-free, sound on the
pipeline, and the only mechanism measured to beat the per-query caps
there; the budget scaling of their constant is the sharpest open
problem of the session (Open Problem~\ref{op:adjacency}).
\end{remark}

\section{Budget reconciliation}\label{sec:budget}

Theorem~\ref{thm:chi} (Theorem 6.1 of \cite{Kra26}) quantifies over
\emph{budgeted} trees: $(d,e')$-trees with $e'=d\log k$ queries (the
printed budget is $e'=O(\log n)+O(d\log n)$; we keep the shorthand
$d\log k$).  Every attacker constructed in this paper is over that
budget in the intended regime $e'=\mathrm{poly}(\log n)\ll n$:

\begin{center}
\small
\begin{tabular}{lll}
\hline
attacker & queries & error (coin channel; chain: measured)\\
\hline
two-phase tree (Theorem~\ref{thm:twophase}) & $\asymp\,2n+1$ &
$5\cdot 2^{-(2d+1)}$\\
rc parity tree (Proposition~\ref{prop:rcparity}) & $2n+1$ &
$(6d+2)\,2^{-6d}$\\
full-scan counting tree (Theorem~\ref{thm:floor}) & $n(n+1)$ &
$2^{-\Theta(d^2)}$\\
confirmation chain (Lem.~\ref{lem:adjacency}) & $B$ (single variables) &
$\approx e^{-B/400}$, decaying in $B$\\
\hline
\end{tabular}
\end{center}

The first three errors are coin-channel closed forms (for the counting
value the $2^{-\Theta(d^2)}$ order is kernel-faithful,
Remark~\ref{rem:parityfloor}); against the literal pipeline the first
two are dead outright (Corollary~\ref{cor:closed}).  The chain's error
is measured on both channels (Remark~\ref{rem:chainmeasured}) and is
the first to decay in budget rather than in $d$.  The two-phase
witness alone exceeds the budget $d\log k$ whenever
$d<2n^{1-\delta}$ at $k=2^{n^{\delta}}$, \emph{a fortiori} the counting
tree.  Counting certificates are inert within budget: a single
variable answers $1$ with probability $m+f/2=O(1/n)$, so $e$
single-variable queries elicit an expected number of $1$s of order
$e/n$, and a line-count certificate (two $1$s in one row or column)
does not occur until $e\gtrsim n$, far above $d\log k$.  Within budget
the certificate mechanism of Lemma~\ref{lem:certificates} cannot fire,
the cap of Theorem~\ref{thm:B} governs (its slack $O(e/n)$ is $o(1)$
there), and the chi-hypothesis survives every known attack at the
printed budget.  The one attacker whose measured error decays in
budget, the confirmation chain, needs $\approx c^{-1}\log(1/\gamma)$
queries with $c\approx1/400$ per query at the measured points; whether
$e'=d\log k$ suffices at the program's parameters is the
adjacency-budget race (Open Problem~\ref{op:adjacency}).  Three
precise statements record where matters stand:

\begin{enumerate}
\item The unbounded degree-$1$ quantifier is settled, by
Theorem~\ref{thm:floor}.  The unbounded quantifier is not the printed
one: Definition~\ref{def:solution} quantifies over $(d,e)$-trees, and
whether any reading of \cite{Kra26} leaves the class unbounded is a
reading question to verify against the source.  Under an unbounded
reading of the solution condition, the counting tree's error
$2^{-\Theta(d^2)}$ falls below the required $\gamma=k^{-O(1)}$
precisely when $4d^2-O(d\log d)>O(\log k)$, which holds throughout the
program's intended regime $d=(\log k)^{O(l)}$ with $l\ge1$ fixed: the
budgeted quantifier of Definition~\ref{def:solution} is essential for
the program's viability at $p=2$.
\item The budgeted floor is open: the optimal error of $e$-query
adaptive degree-$1$ trees as a function of $e$ and, in particular,
whether any $e=d\log k$ tree beats the cap of Theorem~\ref{thm:B}
(Open Problem~\ref{op:floor}).  The confirmation chain answers the
large-$e$ end: it beats every per-query cap for $B\ge300$ at the
measured points, with error decaying like $e^{-B/400}$
(Section~\ref{sec:adjacency}); the open span is the scaling in
between (Open Problems~\ref{op:floor} and~\ref{op:adjacency}).
\item The chi-hypothesis of Theorem~\ref{thm:chi} remains open in
exactly its printed budgeted form (Open Problem~\ref{op:main}).
\end{enumerate}

\section{Consequences for the chi-task}\label{sec:consequences}

Collecting the proved caps, at $p=2$, on the pipeline
$\Omega(n,d)$:
\begin{enumerate}
\item \emph{Non-adaptive variable/monomial-class trees}
(Proposition~\ref{prop:nonadaptall}): success
$\le(\gamma s+1)f+o(1)$; the chi-hypothesis is safe within the class
up to $d=\Theta_\gamma(n^{2/3}(\log k)^{-1/3})$ at budget $s=d\log k$.
The all-queries generality previously claimed for this cap is
retracted (Remark~\ref{rem:retraction}).
\item \emph{All adaptive single-variable trees within budget}
(Theorem~\ref{thm:B}, Remark~\ref{rem:Bscope}): success
$\le\max(q,f)+o(1)$; the hypothesis is
safe with large margin for all $2d\le\sqrt{2n}$, i.e.\ for
$k=2^{n^{\delta}}$ with $\delta\lesssim 1/(2O(l))$ under the program's
parameter connection $d=(\log k)^{O(l)}$.  Without budget the same
class is settled exactly at the Bayes floor $2^{-\Theta(d^2)}$ in the
coin model (Theorem~\ref{thm:floor}); the floor's order is
kernel-faithful and its constant is not
(Remark~\ref{rem:parityfloor}).
\item \emph{Single-pass degree-$2$ product tests}
(Theorem~\ref{thm:and}): no per-hit posterior lift; sharp, and
measured on the exact pipeline (Remark~\ref{rem:andpipeline}).  The
tree-level lift goes through adjacency and needs only single
variables (Section~\ref{sec:adjacency}).
\item \emph{Certification-style attacks and the certificate
trichotomy:} in the coin model
the two-phase tree errs $5\cdot 2^{-(2d+1)}$ under the literal query
semantics (Theorem~\ref{thm:twophase}); it is dominated at equal
$O(n)$ cost by the rc parity tree (Proposition~\ref{prop:rcparity})
and, at $n(n+1)$ queries, by the counting tree's Bayes-optimal
$2^{-\Theta(d^2)}$ (Theorem~\ref{thm:floor}).  Against the literal
pipeline the parity mechanism is dead outright
(Corollary~\ref{cor:closed}), the counting mechanism survives with a
parity-locked hard class (Remark~\ref{rem:parityfloor}), and the
adjacency mechanism is alive and measured: the confirmation chain
reaches $0.9417$ on the exact pipeline at $(32,2)$, budget $1000$
(Section~\ref{sec:adjacency}).  Every unbounded attacker is over
budget, counting certificates are inert below $e\asymp n$
(Section~\ref{sec:budget}), and the chain's reach at the printed
budget is exactly the adjacency-budget race
(Open Problem~\ref{op:adjacency}).  Hence the candidate $\Omega(n,2)$ remains
consistent with all known \emph{budgeted} attackers at $p=2$, and both
the solution condition and the chi-hypothesis remain open and
plausible in their printed budgeted forms.
\item \emph{Trivial multiplicative tests do not refute the candidate:}
an adaptive degree-$2$ tree built from random monomial triples detects
violations of multiplicativity of $\omega=L\circ\rho$ outside the free
region with measured probability $5\times10^{-4}$ over the full
pipeline ($4000$ samples), against $0.37$ for a canonical computation
that drops the restriction layer: $\omega$ is automatically
multiplicative outside the tiny free region, whose all-free-triple
density is $\Theta((d/n)^4)\to0$.
\end{enumerate}

The genuinely open classes are: budgeted degree-$1$ trees at
intermediate budgets (the budgeted floor, with the confirmation chain
as its strongest measured witness), and adaptive trees of degree
$\ge2$ beyond the adjacency mechanism; Section~\ref{sec:open}
formalizes the remaining problems.

\subsection{Theorem-faithfulness}\label{sec:faithfulness}

The coin channel $\mathsf{Coin}(n,d)$ and the literal pipeline agree
\emph{exactly} (total variation $0$) on the marginal and joint
single-variable laws of every row-incomplete query set, and disagree
on a fully queried free row (conditional total variation $\tfrac12$;
parity-locking, Theorem~\ref{thm:channel}(iv)) and at degree $2$
(Remark~\ref{rem:andpipeline}).  Table~\ref{tab:faithfulness} records
the status of each result; ``kernel-faithful'' means the proved
statement holds verbatim for Definition~\ref{def:pipeline}'s pipeline.

\begin{table}[t]
\centering
\footnotesize
\begin{tabular}{lll}
\hline
result & status & why\\
\hline
coin equivalence (Cor.~\ref{cor:coin}) & kernel-faithful & exact off
$E_{\mathrm{full}}$ (TV $=0$); $\tfrac12$ on full rows\\
channel law (Thm.~\ref{thm:channel}) & kernel-faithful & it \emph{is} the
kernel law (verified)\\
completion (Lem.~\ref{lem:completion}) & kernel-faithful & the zeroing
completion, all $d$\\
Prop.~\ref{prop:fixed} (A), \ref{prop:nonadaptsingle} (C) & kernel-faithful &
pair symmetry; single variables\\
Prop.~\ref{prop:nonadaptall} (D, variable/monomial class) & kernel-faithful &
Lemma~\ref{lem:factor} $+$ depletion\\
Thm.~\ref{thm:B}, Bayes caps (Lem.~\ref{lem:bayes}) & kernel-faithful &
single variables; TV $=0$ off $E_{\mathrm{full}}$\\
counting certificates (Lem.~\ref{lem:certificates}) & kernel-faithful &
determined answers; odd-count law corrected\\
Thm.~\ref{thm:floor}'s \emph{order} $2^{-\Theta(d^2)}$ & kernel-faithful &
Rem.~\ref{rem:parityfloor}\\
adjacency certificate (Lem.~\ref{lem:adjacency}) & kernel-faithful &
determined answers only\\
\hline
two-phase law (Thm.~\ref{thm:twophase}) & coin-model-only & parity mechanism
dead (Cor.~\ref{cor:closed})\\
rc parity tree (Prop.~\ref{prop:rcparity}) & coin-model-only & row sums answer
$0$ (Prop.~\ref{prop:rigidity})\\
Thm.~\ref{thm:floor}'s \emph{constant} & coin-model-only & hard class
parity-locked (Rem.~\ref{rem:parityfloor})\\
AND no-lift identity (Thm.~\ref{thm:and}) & coin-model-only$^{\dagger}$ &
product semantics; transfer unproved\\
\hline
\end{tabular}
\caption{Theorem-faithfulness against the literal pipeline of
Definition~\ref{def:pipeline}.  ``Kernel-faithful'': the statement holds
verbatim for the pipeline (the coin channel agrees with it exactly on the
query sets the result uses).  ``Coin-model-only'': the statement is proved
in $\mathsf{Coin}(n,d)$ and does not transfer verbatim; for
Theorem~\ref{thm:and} the pipeline transfer is measured-consistent but
unproved.  Results above the second rule are used in their
kernel-faithful form throughout Sections~\ref{sec:budget}
and~\ref{sec:open}.}
\label{tab:faithfulness}
\end{table}

\noindent$^{\dagger}$Measured-consistent: true monomial queries on the exact
pipeline land on the cap, $0.2629$ $[0.2129,0.3199]$ against
$q=0.2632$ (Remark~\ref{rem:andpipeline}).

\section{Open problems}\label{sec:open}

\begin{openproblem}[the chi-hypothesis, printed budgeted form]
\label{op:main}
Is there a constant $C$ such that every $(d,e')$-tree with
$e'=d\log k$ (printed: $e'=O(\log n)+O(d\log n)$) and queries of
degree $\le d$ has, on the pipeline
$\Omega(n,d)$ at $p=2$ with $d=(\log k)^{O(l)}\le\sqrt{n}$,
\[
\Prob[\text{failure}]\;\ge\;k^{-C}\;?
\]
Equivalently: no shallow low-degree adaptive tree certifies freeness
with probability $1-k^{-O(1)}$.  An affirmative answer, together with
Theorem~\ref{thm:chi}, yields $k$-step lower bounds for
$\AC^0[l](\mathrm{MOD}_2)$-Frege refutations of $\PHPn$ with
$k=2^{n^{\delta}}$, $\delta\approx1/(2O(l))$: the first
super-polynomial $\AC[2]$-Frege lower bounds, on the pigeonhole
principle.  A negative answer would refute $\Omega(n,d)$ as a
pseudo-solution at $p=2$ and redirect the program.  Proved here: the
affirmative answer for non-adaptive variable/monomial-class trees
(Proposition~\ref{prop:nonadaptall}) and for adaptive single-variable
trees at budget (Theorem~\ref{thm:B}); without budget, degree $1$ is
settled exactly at the Bayes floor $2^{-\Theta(d^2)}$
(Theorem~\ref{thm:floor}), which stays above $k^{-O(1)}$ only while
$d^2\lesssim\log k$ --- one more reason the budgeted quantifier is the
right one.  The residual classes: budgeted degree-$1$ trees at
intermediate budgets (Open Problem~\ref{op:floor}), now dominated in
the measured range by the confirmation chain, whose error decays in
budget like $e^{-B/400}$ on the pipeline itself
(Section~\ref{sec:adjacency}; its budget scaling is
Open Problem~\ref{op:adjacency}), and adaptive trees of degree $\ge2$
beyond the adjacency mechanism (Open Problem~\ref{op:chaining}).
\end{openproblem}

\begin{openproblem}[chaining beyond the per-hit posterior]
\label{op:chaining}
The per-hit no-lift is sharp: every single-pass test has posterior
exactly $q$ (Theorem~\ref{thm:and}), measured on the exact pipeline as
well (Remark~\ref{rem:andpipeline}).  At the tree level the picture
inverted: adaptive chaining \emph{does} beat the per-query posterior
cap --- not by raising any single answer's posterior, but by chaining
a hit with an adjacent answer into a certain-free event whose
probability grows exponentially with budget (Section~\ref{sec:adjacency}:
$0.9317$ against the cap $0.2550$ at $(32,2)$, budget $1000$, disjoint
$99\%$ intervals).  The earlier ``every tested chaining strategy decays
as $n^{-1.8}$ to $n^{-1.9}$'' reading is superseded; the probe behind
it never implemented its product-test branch (a measurement bug
recorded in the corpus).  What remains open: (a) the budget scaling of
the adjacency constant (Open Problem~\ref{op:adjacency}); (b) whether
degree-$2$ chains add anything over the degree-$1$ confirmation chain
--- budget-honest AND-scans are \emph{worse} than single variables,
and AND-split chains certify far more rarely ($4.00\%$ versus
$91.83\%$ certificate rate at $(32,2)$, budget $1000$); (c) a cap
theorem for the adaptive degree-$1$ class at intermediate budgets,
between Theorem~\ref{thm:B}'s $e=o(n)$ regime and the unbounded floor
(Theorem~\ref{thm:floor}).
\end{openproblem}

\begin{openproblem}[the budgeted certification floor]\label{op:floor}
What is the optimal error $\mathrm{err}(e)$ of adaptive degree-$1$
trees with $e$ queries, as a function of $e$?  Unbounded, the answer
is exact: $\mathrm{err}^*=2^{-\Theta(d^2)}$, attained by the
non-adaptive counting tree (Theorem~\ref{thm:floor}); the earlier
conjecture of a $2^{-\Theta(d)}$ floor, with the two-phase tree as its
extremal witness, is false.  The line-count certificates that attain
the floor are inert below $e\asymp n$ (Section~\ref{sec:budget}), and
Theorem~\ref{thm:B} caps the class at $\max(q,f)(1+o(1))$ for
$e=o(n)$; the question is whether the truth interpolates or some
$e=d\log k$ tree beats Theorem~\ref{thm:B}'s cap, which is the content
Theorem~\ref{thm:chi}'s hypothesis needs at $p=2$.  An affirmative
budgeted floor would combine with Open Problem~\ref{op:main} to
complete the $p=2$ core of the $\AC[2]$-Frege program.  Three
sub-questions: (a) the exact Bayes program at degree $2$ ---
Corollary~\ref{cor:fullscan}'s reduction dies there (a monomial answer
is not an XOR of single answers), so no degree-$\ge2$ analogue of
Theorem~\ref{thm:floor} is known; (b) against the literal pipeline,
the best statistical freeness signal obtainable from aggregate
(full-row or higher-degree) queries within budget, the quantitative
remnant of Corollary~\ref{cor:closed}; (c) the \emph{parity-locked
floor constant}: recompute Theorem~\ref{thm:floor}'s constant for the
parity-locked answer table, in which every free row carries an odd
count of $1$s ($\Prob[\text{count}=1]=2d/2^{2d-1}$, twice the coin
model's) and no free row is ever all-zero --- the $2^{-\Theta(d^2)}$
order survives (Remark~\ref{rem:parityfloor}), the exact constant does
not, and it is open.
\end{openproblem}

\begin{openproblem}[the adjacency-budget race]\label{op:adjacency}
The confirmation chain's certificate rate is
$\approx1-e^{-cB}$ at budget $B$, with $c\approx1/400$ per query
measured at both $(32,2)$ and $(64,4)$ (Remark~\ref{rem:chainmeasured});
its error on the exact pipeline decays like $e^{-B/400}$.  Determine
$c(n,d)$.  The two measured points have nearly equal free densities
$f$ and hit rates, so the scaling of $c$ in $(n,d)$ is untested; the
free-density heuristic of the chaining study gives
$c\sim d^2/n^2$ to $d^3/n^3$ (each read hits a free pair with
probability $\approx f$, and a free hit then certifies within a few
neighbor reads), under which the budget needed to push the error below
$\gamma$ scales as $(n^3/d^3)\log(1/\gamma)$: polynomial in $n$,
polylogarithmic in $1/\gamma$.  The stakes are exact: if
$e'=d\log k$ queries suffice at the program's parameters
$d=(\log k)^{O(l)}$, the printed chi-hypothesis of
Theorem~\ref{thm:chi} is \textbf{false} and the $p=2$ route via
Theorem~\ref{thm:chi} closes; if $c$ decays polynomially in $n$, the
hypothesis survives with the confirmation chain as its strongest known
witness.  This is the sharpest open problem of the session: the
budgeted certification floor must cover a trichotomy of certificate
mechanisms --- parity (dead on the pipeline,
Corollary~\ref{cor:closed}), counting (inert within polylog budget,
Section~\ref{sec:budget}), and adjacency (alive; this problem).
\end{openproblem}

\begin{openproblem}[the odd-characteristic analogue]\label{op:oddprime}
Fix an odd prime $p$ (any $p\ge3$).  The reduction chain of
\cite{Kra26} is characteristic-uniform (Theorems 2.2, 3.2, 3.3,
Definition 4.3, and Theorem 6.1 there are all stated for an arbitrary
prime), and the free-answer alphabet grows to all of $\F_p$: killed
lines still answer determinedly, free single variables answer uniform
$p$-ary dies, so a fraction $\frac{p-2}{p}$ of free single-variable
answers are self-certifying and the single-variable lane is open at
odd $p$ in a way it is not at $p=2$.  Determine whether, for every
$(d,e')$-tree $T'$ with $d=(\log k)^{O(l)}$ and
$e'=O(\log n)+O(d\log n)$,
\[
\Prob_{\rho,P}\bigl[\chi(P,\rho)=1\bigr]\;\ge\;k^{-O(1)},
\]
the hypothesis of Theorem~\ref{thm:chi} at characteristic $p$.  A
positive answer at any $k(n)\ge n^{\omega(1)}$ yields, via the
Extended Nullstellensatz equivalence, $k$-step lower bounds for
$\AC^0[p]$-Frege refutations of $\PHPn$: the first super-polynomial
lower bound for the full $\AC^0[p]$-Frege system at that
characteristic, and the first lower bound of any kind for odd-$p$
$\AC^0[p]$-Frege proofs of the pigeonhole principle; the odd-$p$
parity-resolution lane is strictly less crowded (no dag-like
Res$(\mathrm{lin}_p)$ pigeonhole bound of any kind exists for odd
$p$, only tree-like and fragment results).  A negative answer (a tree
with error $<k^{-O(1)}$) refutes $\Omega(n,d)$ at that $p$ and
redirects the program, exactly as at $p=2$.  Caveat: the $p$-ary
two-phase tree's proposed error $p^{-(2d+1)}$ is flagged
\emph{analysis} in the corpus; it inherits the phase-$2$ coin-branch
issue corrected in Theorem~\ref{thm:twophase} (the certified row's
sum being nonzero does not prevent the all-zero configuration), and
its literal error needs the same recomputation before any use.
\end{openproblem}

\section*{Acknowledgements}
The corpus owner supplied the research environment and the verification
scripts.  All computational claims are reproducible from the session's
exact channel samplers, from the registered run of
\texttt{cert\_floor\_check.py} (exact enumeration, Bayes audit, and
Monte Carlo, $31$\,s), from the chaining study's registered run
(\texttt{chi\_and\_chain.py}: $600$ simulations per cell, Wilson
$99\%$ intervals, asserted certifications, $52{,}800$ simulations,
seed $20261003$), and from the kernel-level verification
(\texttt{kernel\_structure.py}: $18$ checks and the one-kernel,
parity-locking, and degree-$2$ classifications, $100{,}000$ design
samples per law).

\begin{thebibliography}{99}

\bibitem{Ajtai94}
M.~Ajtai.
\newblock The complexity of the pigeonhole principle.
\newblock \emph{Combinatorica} 14(4):417--433, 1994.

\bibitem{BIKPPS97}
S.~R.~Buss, R.~Impagliazzo, J.~Kraj\'i\v{c}ek, P.~Pudl\'ak,
A.~A.~Razborov, and J.~Sgall.
\newblock Proof complexity in algebraic proof systems.
\newblock \emph{Journal of the ACM} 44(1):166--191, 1997.

\bibitem{CookReckhow79}
S.~A.~Cook and R.~A.~Reckhow.
\newblock The relative efficiency of propositional proof systems.
\newblock \emph{Journal of Symbolic Logic} 44(1):36--50, 1979.

\bibitem{DubhashiRanjan98}
D.~Dubhashi and D.~Ranjan.
\newblock Balls and bins: a study in negative dependence.
\newblock \emph{Random Structures \& Algorithms} 13(2):99--124, 1998.

\bibitem{Haken85}
A.~Haken.
\newblock The intractability of resolution.
\newblock \emph{Theoretical Computer Science} 39:297--308, 1985.

\bibitem{JoagDevProschan83}
K.~Joag-Dev and F.~Proschan.
\newblock Negative association of random variables, with applications.
\newblock \emph{The Annals of Statistics} 11(1):286--295, 1983.

\bibitem{Kra19}
J.~Kraj\'i\v{c}ek.
\newblock \emph{Proof Complexity}.
\newblock Encyclopedia of Mathematics and its Applications 170, Cambridge
University Press, 2019.

\bibitem{Kra26}
J.~Kraj\'i\v{c}ek.
\newblock Pseudo-solutions of polynomial systems and the lower bound
problem for $\AC^0[p]$-Frege systems.
\newblock arXiv:2609.35927, September 2026.

\bibitem{randomcnf26}
Super-linear lower bounds for random $3$-CNFs in bounded-depth Frege.
\newblock arXiv:2403.02275, 2024--2026.

\bibitem{resxor26}
An exponential lower bound for the bit pigeonhole principle in
unrestricted dag-like resolution over parities.
\newblock arXiv:2609.23015, 2026; Lean~4 formalization.

\bibitem{Smolensky87}
R.~Smolensky.
\newblock Algebraic methods in the theory of lower bounds for Boolean
circuit complexity.
\newblock In \emph{Proceedings of the 19th Annual ACM Symposium on
Theory of Computing (STOC)}, pages 77--82, 1987.

\end{thebibliography}

\end{document}
